%==============================================================================
% Unconditional Primality Certificate — 29 998-digit prime
% Parametric family  p = 3m(m+1)+1,  m = 2^a·3^b − 1
% Independent numerical verification — May 2026
%==============================================================================
%
% Stand-alone usage (requires paper_pocklington.tex preamble, or use the
% minimal preamble at the bottom of this file):
%
%   %==============================================================================
% Unconditional Primality Certificate — 29 998-digit prime
% Parametric family  p = 3m(m+1)+1,  m = 2^a·3^b − 1
% Independent numerical verification — May 2026
%==============================================================================
%
% Stand-alone usage (requires paper_pocklington.tex preamble, or use the
% minimal preamble at the bottom of this file):
%
%   %==============================================================================
% Unconditional Primality Certificate — 29 998-digit prime
% Parametric family  p = 3m(m+1)+1,  m = 2^a·3^b − 1
% Independent numerical verification — May 2026
%==============================================================================
%
% Stand-alone usage (requires paper_pocklington.tex preamble, or use the
% minimal preamble at the bottom of this file):
%
%   %==============================================================================
% Unconditional Primality Certificate — 29 998-digit prime
% Parametric family  p = 3m(m+1)+1,  m = 2^a·3^b − 1
% Independent numerical verification — May 2026
%==============================================================================
%
% Stand-alone usage (requires paper_pocklington.tex preamble, or use the
% minimal preamble at the bottom of this file):
%
%   \input{certificate_29998}          % inside the main document
%
% Self-contained compilation:
%   pdflatex certificate_29998.tex
%
%==============================================================================

%------------------------------------------------------------------------------
\section{Unconditional Certificate for the 29\,998-Digit Prime}
\label{sec:cert29998}
%------------------------------------------------------------------------------

This section provides the complete, independently verified Pocklington–Lehmer
primality certificate for the largest prime currently produced by
\textsc{PrimeQuest}.

%------------------------------------------------------------------------------
\subsection{Parameters and Structure}
%------------------------------------------------------------------------------

\begin{definition}[29\,998-digit PrimeQuest instance]\label{def:instance29998}
Set
\[
  a = 19\,435,\qquad b = 19\,173,
\]
and define
\[
  m \;=\; 2^{19435}\cdot 3^{19173} - 1
    \qquad\bigl(14\,999 \text{ decimal digits}\bigr),
\]
\[
  p \;=\; 3m(m+1)+1
    \qquad\bigl(29\,998 \text{ decimal digits}\bigr).
\]
\end{definition}

The ratio $b/a = 19\,173/19\,435 \approx 0.9865$ is close to the
theoretical optimum
\[
  \rho^* = \frac{\log_{10} 2}{\log_{10} 3} \approx 0.6309
\]
in the sense that it minimises the digit count for a prescribed $a$;
the value used here targets exactly $29\,998$ digits and was located by
the \textsc{PrimeQuest v4} algorithm (see Section~\ref{sec:algo}).

\medskip

\noindent\textbf{Identification.}\quad
The first and last decimal digits of $p$ are:
\begin{equation}\label{eq:p29998digits}
  p \;=\; 1602087359509060348500037851549319056715\cdots
         16930852697403293697.
\end{equation}

%------------------------------------------------------------------------------
\subsection{Factorisation of $p-1$}
%------------------------------------------------------------------------------

\begin{rigorousbox}
\begin{lemma}[Factorisation of $p-1$]\label{lem:factorisation29998}
With $a, b, m, p$ as in Definition~\ref{def:instance29998},
\[
  p - 1 \;=\; \underbrace{2^{19435}\cdot 3^{19174}}_{F}
              \;\cdot\; m,
\]
where $F := 2^{19435}\cdot 3^{19174}$ has $15\,000$ decimal digits and
satisfies $F > \sqrt{p}$.
\end{lemma}
\end{rigorousbox}

\begin{proof}
Since $m + 1 = 2^{19435}\cdot 3^{19173}$,
\[
  p - 1 \;=\; 3m(m+1) \;=\; 3m\cdot 2^{19435}\cdot 3^{19173}
           \;=\; 2^{19435}\cdot 3^{19174}\cdot m.
\]
For the bound $F > \sqrt{p}$, note that $p < 3(m+1)^2$, so
$\sqrt{p} < \sqrt{3}\,(m+1) < 3(m+1) = 2^{19435}\cdot 3^{19174} = F$.

\noindent Numerically, comparing base-$2$ logarithms:
\[
  2\log_2 F
  \;=\; 2\bigl(19435 + 19174\log_2 3\bigr)
  \;\approx\; 99\,650.14,
\]
\[
  \log_2 p \;\approx\; 99\,648.56,
\]
confirming $F^2 > p$ with a margin of $\approx 1.58$ bits.
\end{proof}

%------------------------------------------------------------------------------
\subsection{The Pocklington--Lehmer Certificate}
%------------------------------------------------------------------------------

Theorem~\ref{thm:Pocklington} (recalled from Section~\ref{sec:family})
requires, for each prime $q \mid F$ — here $q \in \{2, 3\}$ — an integer
$w_q$ satisfying
\[
  w_q^{p-1} \equiv 1 \pmod{p}
  \qquad\text{and}\qquad
  \gcd\!\bigl(w_q^{(p-1)/q} - 1,\; p\bigr) = 1.
\]

\begin{resultbox}[title={\textbf{Certificate} — 29\,998-digit prime}]
\[
  p \;=\; 3m(m+1)+1,\qquad m = 2^{19435}\cdot 3^{19173}-1.
\]

\medskip
\noindent\textbf{Factored part of $p-1$:}
$\;F = 2^{19435}\cdot 3^{19174},\quad F > \sqrt{p}.\quad$\checkmark

\medskip
\renewcommand{\arraystretch}{1.4}
\begin{tabular}{@{}c l l@{}}
\toprule
Prime $q$ & Witness $w_q$ & Conditions verified \\
\midrule
$2$ & $w_2 = 5$ &
  $5^{p-1}\equiv 1\pmod{p}$\quad and\quad
  $\gcd(5^{(p-1)/2}-1,\,p)=1$ \\
$3$ & $w_3 = 7$ &
  $7^{p-1}\equiv 1\pmod{p}$\quad and\quad
  $\gcd(7^{(p-1)/3}-1,\,p)=1$ \\
\bottomrule
\end{tabular}

\medskip
\noindent\textbf{Conclusion:} $p$ is prime (deterministic proof). \checkmark
\end{resultbox}

%------------------------------------------------------------------------------
\subsection{Independent Numerical Verification}
%------------------------------------------------------------------------------

All four conditions of the certificate were verified independently using
\texttt{Python 3.12} with its built-in arbitrary-precision integer
arithmetic (GMP backend). The script performs no probabilistic step;
every operation is exact modular arithmetic.

\begin{rigorousbox}[title={\textbf{Verification record} — independent computation}]
\textbf{Environment:} Python 3.12, \texttt{sys.set\_int\_max\_str\_digits(40000)};
modular exponentiations via \texttt{pow(base, exp, mod)}.\\[4pt]

\textbf{Step 1 — Reconstructing $p$.}\quad
Computed $m = 2^{19435}\cdot 3^{19173}-1$ (14\,999 digits) and
$p = 3m(m+1)+1$ (29\,998 digits) from scratch.
First 40 digits of $p$: \texttt{1602087359509060348500037851549319056715};
last 20 digits: \texttt{16930852697403293697}.
Both match the stored certificate exactly.\\[4pt]

\textbf{Step 2 — Pocklington bound.}
\[
  2\log_2 F = 99\,650.1420 \;>\; \log_2 p = 99\,648.5570. \checkmark
\]

\textbf{Step 3 — Witness $w_2 = 5$ for $q = 2$.}
\begin{align*}
  5^{p-1} &\equiv 1 \pmod{p}, \checkmark \\
  \gcd(5^{(p-1)/2}-1,\,p) &= 1. \checkmark
\end{align*}

\textbf{Step 4 — Witness $w_3 = 7$ for $q = 3$.}
\begin{align*}
  7^{p-1} &\equiv 1 \pmod{p}, \checkmark \\
  \gcd(7^{(p-1)/3}-1,\,p) &= 1. \checkmark
\end{align*}

\textbf{Wall-clock time:} $6\,945$ s $\approx 1$ h $56$ min
(single core, Python GMP, no multi-threading).\\[4pt]

\textbf{Conclusion:} all four Pocklington–Lehmer conditions are satisfied.
The number $p$ is \emph{unconditionally prime}.
\end{rigorousbox}

%------------------------------------------------------------------------------
\subsection{Comparison with the Original Computation}
%------------------------------------------------------------------------------

\begin{table}[h]
\centering
\caption{Original search vs.\ independent certificate verification
         for the 29\,998-digit prime.}
\label{tab:comparison29998}
\renewcommand{\arraystretch}{1.3}
\begin{tabular}{@{}lll@{}}
\toprule
 & \textbf{PrimeQuest v4 (search)} & \textbf{Independent verification} \\
\midrule
Purpose        & Find a prime candidate          & Check the certificate \\
Machine        & ARM64, 9 cores (Snapdragon X)   & Single core (Python GMP) \\
Wall-clock     & $11\,146$ s ($3$ h $06$ min)    & $6\,945$ s ($1$ h $56$ min) \\
MR tests run   & $286$ (probabilistic pre-filter) & — (no probabilistic step) \\
Result         & Certificate produced             & Certificate confirmed \\
\bottomrule
\end{tabular}
\end{table}

The verification is $\approx 1.6\times$ faster than the original search
because it executes exactly four modular exponentiations on the known
candidate, bypassing the sieve, Theorem-2 filter, and Miller–Rabin
pre-screening required during the search phase.

%------------------------------------------------------------------------------
\subsection{Record in Context}
%------------------------------------------------------------------------------

Table~\ref{tab:allcerts} (Section~\ref{sec:results}) lists all four
\textsc{PrimeQuest} certificates. The 29\,998-digit instance is currently
the largest unconditional Pocklington–Lehmer certificate in the family
$p = 3m(m+1)+1$, $m = 2^a\cdot 3^b - 1$, produced on a single consumer
laptop without distributed computing or specialised hardware.

%==============================================================================
% MINIMAL STAND-ALONE PREAMBLE (uncomment for self-contained compilation)
%==============================================================================
% \iffalse   % <-- remove this line to enable stand-alone mode
% \documentclass[11pt,a4paper]{article}
% \usepackage[utf8]{inputenc}
% \usepackage[T1]{fontenc}
% \usepackage[english]{babel}
% \usepackage{amsmath,amssymb,amsthm,mathtools,booktabs}
% \usepackage{geometry}\geometry{margin=2.5cm}
% \usepackage[dvipsnames]{xcolor}
% \usepackage{tcolorbox}
% \tcbuselibrary{skins,breakable}
% \newtcolorbox{rigorousbox}[1][]{colback=Green!8,colframe=Green!60!black,
%   arc=1mm,title=\textbf{Rigorous},fonttitle=\bfseries,breakable,#1}
% \newtcolorbox{resultbox}[1][]{colback=RoyalBlue!5,colframe=RoyalBlue!70!black,
%   arc=1mm,title=\textbf{Computational Result},fonttitle=\bfseries,breakable,#1}
% \newtheorem{theorem}{Theorem}[section]
% \newtheorem{lemma}[theorem]{Lemma}
% \theoremstyle{definition}
% \newtheorem{definition}[theorem]{Definition}
% \begin{document}
% \input{certificate_29998}
% \end{document}
% \fi
          % inside the main document
%
% Self-contained compilation:
%   pdflatex certificate_29998.tex
%
%==============================================================================

%------------------------------------------------------------------------------
\section{Unconditional Certificate for the 29\,998-Digit Prime}
\label{sec:cert29998}
%------------------------------------------------------------------------------

This section provides the complete, independently verified Pocklington–Lehmer
primality certificate for the largest prime currently produced by
\textsc{PrimeQuest}.

%------------------------------------------------------------------------------
\subsection{Parameters and Structure}
%------------------------------------------------------------------------------

\begin{definition}[29\,998-digit PrimeQuest instance]\label{def:instance29998}
Set
\[
  a = 19\,435,\qquad b = 19\,173,
\]
and define
\[
  m \;=\; 2^{19435}\cdot 3^{19173} - 1
    \qquad\bigl(14\,999 \text{ decimal digits}\bigr),
\]
\[
  p \;=\; 3m(m+1)+1
    \qquad\bigl(29\,998 \text{ decimal digits}\bigr).
\]
\end{definition}

The ratio $b/a = 19\,173/19\,435 \approx 0.9865$ is close to the
theoretical optimum
\[
  \rho^* = \frac{\log_{10} 2}{\log_{10} 3} \approx 0.6309
\]
in the sense that it minimises the digit count for a prescribed $a$;
the value used here targets exactly $29\,998$ digits and was located by
the \textsc{PrimeQuest v4} algorithm (see Section~\ref{sec:algo}).

\medskip

\noindent\textbf{Identification.}\quad
The first and last decimal digits of $p$ are:
\begin{equation}\label{eq:p29998digits}
  p \;=\; 1602087359509060348500037851549319056715\cdots
         16930852697403293697.
\end{equation}

%------------------------------------------------------------------------------
\subsection{Factorisation of $p-1$}
%------------------------------------------------------------------------------

\begin{rigorousbox}
\begin{lemma}[Factorisation of $p-1$]\label{lem:factorisation29998}
With $a, b, m, p$ as in Definition~\ref{def:instance29998},
\[
  p - 1 \;=\; \underbrace{2^{19435}\cdot 3^{19174}}_{F}
              \;\cdot\; m,
\]
where $F := 2^{19435}\cdot 3^{19174}$ has $15\,000$ decimal digits and
satisfies $F > \sqrt{p}$.
\end{lemma}
\end{rigorousbox}

\begin{proof}
Since $m + 1 = 2^{19435}\cdot 3^{19173}$,
\[
  p - 1 \;=\; 3m(m+1) \;=\; 3m\cdot 2^{19435}\cdot 3^{19173}
           \;=\; 2^{19435}\cdot 3^{19174}\cdot m.
\]
For the bound $F > \sqrt{p}$, note that $p < 3(m+1)^2$, so
$\sqrt{p} < \sqrt{3}\,(m+1) < 3(m+1) = 2^{19435}\cdot 3^{19174} = F$.

\noindent Numerically, comparing base-$2$ logarithms:
\[
  2\log_2 F
  \;=\; 2\bigl(19435 + 19174\log_2 3\bigr)
  \;\approx\; 99\,650.14,
\]
\[
  \log_2 p \;\approx\; 99\,648.56,
\]
confirming $F^2 > p$ with a margin of $\approx 1.58$ bits.
\end{proof}

%------------------------------------------------------------------------------
\subsection{The Pocklington--Lehmer Certificate}
%------------------------------------------------------------------------------

Theorem~\ref{thm:Pocklington} (recalled from Section~\ref{sec:family})
requires, for each prime $q \mid F$ — here $q \in \{2, 3\}$ — an integer
$w_q$ satisfying
\[
  w_q^{p-1} \equiv 1 \pmod{p}
  \qquad\text{and}\qquad
  \gcd\!\bigl(w_q^{(p-1)/q} - 1,\; p\bigr) = 1.
\]

\begin{resultbox}[title={\textbf{Certificate} — 29\,998-digit prime}]
\[
  p \;=\; 3m(m+1)+1,\qquad m = 2^{19435}\cdot 3^{19173}-1.
\]

\medskip
\noindent\textbf{Factored part of $p-1$:}
$\;F = 2^{19435}\cdot 3^{19174},\quad F > \sqrt{p}.\quad$\checkmark

\medskip
\renewcommand{\arraystretch}{1.4}
\begin{tabular}{@{}c l l@{}}
\toprule
Prime $q$ & Witness $w_q$ & Conditions verified \\
\midrule
$2$ & $w_2 = 5$ &
  $5^{p-1}\equiv 1\pmod{p}$\quad and\quad
  $\gcd(5^{(p-1)/2}-1,\,p)=1$ \\
$3$ & $w_3 = 7$ &
  $7^{p-1}\equiv 1\pmod{p}$\quad and\quad
  $\gcd(7^{(p-1)/3}-1,\,p)=1$ \\
\bottomrule
\end{tabular}

\medskip
\noindent\textbf{Conclusion:} $p$ is prime (deterministic proof). \checkmark
\end{resultbox}

%------------------------------------------------------------------------------
\subsection{Independent Numerical Verification}
%------------------------------------------------------------------------------

All four conditions of the certificate were verified independently using
\texttt{Python 3.12} with its built-in arbitrary-precision integer
arithmetic (GMP backend). The script performs no probabilistic step;
every operation is exact modular arithmetic.

\begin{rigorousbox}[title={\textbf{Verification record} — independent computation}]
\textbf{Environment:} Python 3.12, \texttt{sys.set\_int\_max\_str\_digits(40000)};
modular exponentiations via \texttt{pow(base, exp, mod)}.\\[4pt]

\textbf{Step 1 — Reconstructing $p$.}\quad
Computed $m = 2^{19435}\cdot 3^{19173}-1$ (14\,999 digits) and
$p = 3m(m+1)+1$ (29\,998 digits) from scratch.
First 40 digits of $p$: \texttt{1602087359509060348500037851549319056715};
last 20 digits: \texttt{16930852697403293697}.
Both match the stored certificate exactly.\\[4pt]

\textbf{Step 2 — Pocklington bound.}
\[
  2\log_2 F = 99\,650.1420 \;>\; \log_2 p = 99\,648.5570. \checkmark
\]

\textbf{Step 3 — Witness $w_2 = 5$ for $q = 2$.}
\begin{align*}
  5^{p-1} &\equiv 1 \pmod{p}, \checkmark \\
  \gcd(5^{(p-1)/2}-1,\,p) &= 1. \checkmark
\end{align*}

\textbf{Step 4 — Witness $w_3 = 7$ for $q = 3$.}
\begin{align*}
  7^{p-1} &\equiv 1 \pmod{p}, \checkmark \\
  \gcd(7^{(p-1)/3}-1,\,p) &= 1. \checkmark
\end{align*}

\textbf{Wall-clock time:} $6\,945$ s $\approx 1$ h $56$ min
(single core, Python GMP, no multi-threading).\\[4pt]

\textbf{Conclusion:} all four Pocklington–Lehmer conditions are satisfied.
The number $p$ is \emph{unconditionally prime}.
\end{rigorousbox}

%------------------------------------------------------------------------------
\subsection{Comparison with the Original Computation}
%------------------------------------------------------------------------------

\begin{table}[h]
\centering
\caption{Original search vs.\ independent certificate verification
         for the 29\,998-digit prime.}
\label{tab:comparison29998}
\renewcommand{\arraystretch}{1.3}
\begin{tabular}{@{}lll@{}}
\toprule
 & \textbf{PrimeQuest v4 (search)} & \textbf{Independent verification} \\
\midrule
Purpose        & Find a prime candidate          & Check the certificate \\
Machine        & ARM64, 9 cores (Snapdragon X)   & Single core (Python GMP) \\
Wall-clock     & $11\,146$ s ($3$ h $06$ min)    & $6\,945$ s ($1$ h $56$ min) \\
MR tests run   & $286$ (probabilistic pre-filter) & — (no probabilistic step) \\
Result         & Certificate produced             & Certificate confirmed \\
\bottomrule
\end{tabular}
\end{table}

The verification is $\approx 1.6\times$ faster than the original search
because it executes exactly four modular exponentiations on the known
candidate, bypassing the sieve, Theorem-2 filter, and Miller–Rabin
pre-screening required during the search phase.

%------------------------------------------------------------------------------
\subsection{Record in Context}
%------------------------------------------------------------------------------

Table~\ref{tab:allcerts} (Section~\ref{sec:results}) lists all four
\textsc{PrimeQuest} certificates. The 29\,998-digit instance is currently
the largest unconditional Pocklington–Lehmer certificate in the family
$p = 3m(m+1)+1$, $m = 2^a\cdot 3^b - 1$, produced on a single consumer
laptop without distributed computing or specialised hardware.

%==============================================================================
% MINIMAL STAND-ALONE PREAMBLE (uncomment for self-contained compilation)
%==============================================================================
% \iffalse   % <-- remove this line to enable stand-alone mode
% \documentclass[11pt,a4paper]{article}
% \usepackage[utf8]{inputenc}
% \usepackage[T1]{fontenc}
% \usepackage[english]{babel}
% \usepackage{amsmath,amssymb,amsthm,mathtools,booktabs}
% \usepackage{geometry}\geometry{margin=2.5cm}
% \usepackage[dvipsnames]{xcolor}
% \usepackage{tcolorbox}
% \tcbuselibrary{skins,breakable}
% \newtcolorbox{rigorousbox}[1][]{colback=Green!8,colframe=Green!60!black,
%   arc=1mm,title=\textbf{Rigorous},fonttitle=\bfseries,breakable,#1}
% \newtcolorbox{resultbox}[1][]{colback=RoyalBlue!5,colframe=RoyalBlue!70!black,
%   arc=1mm,title=\textbf{Computational Result},fonttitle=\bfseries,breakable,#1}
% \newtheorem{theorem}{Theorem}[section]
% \newtheorem{lemma}[theorem]{Lemma}
% \theoremstyle{definition}
% \newtheorem{definition}[theorem]{Definition}
% \begin{document}
% %==============================================================================
% Unconditional Primality Certificate — 29 998-digit prime
% Parametric family  p = 3m(m+1)+1,  m = 2^a·3^b − 1
% Independent numerical verification — May 2026
%==============================================================================
%
% Stand-alone usage (requires paper_pocklington.tex preamble, or use the
% minimal preamble at the bottom of this file):
%
%   \input{certificate_29998}          % inside the main document
%
% Self-contained compilation:
%   pdflatex certificate_29998.tex
%
%==============================================================================

%------------------------------------------------------------------------------
\section{Unconditional Certificate for the 29\,998-Digit Prime}
\label{sec:cert29998}
%------------------------------------------------------------------------------

This section provides the complete, independently verified Pocklington–Lehmer
primality certificate for the largest prime currently produced by
\textsc{PrimeQuest}.

%------------------------------------------------------------------------------
\subsection{Parameters and Structure}
%------------------------------------------------------------------------------

\begin{definition}[29\,998-digit PrimeQuest instance]\label{def:instance29998}
Set
\[
  a = 19\,435,\qquad b = 19\,173,
\]
and define
\[
  m \;=\; 2^{19435}\cdot 3^{19173} - 1
    \qquad\bigl(14\,999 \text{ decimal digits}\bigr),
\]
\[
  p \;=\; 3m(m+1)+1
    \qquad\bigl(29\,998 \text{ decimal digits}\bigr).
\]
\end{definition}

The ratio $b/a = 19\,173/19\,435 \approx 0.9865$ is close to the
theoretical optimum
\[
  \rho^* = \frac{\log_{10} 2}{\log_{10} 3} \approx 0.6309
\]
in the sense that it minimises the digit count for a prescribed $a$;
the value used here targets exactly $29\,998$ digits and was located by
the \textsc{PrimeQuest v4} algorithm (see Section~\ref{sec:algo}).

\medskip

\noindent\textbf{Identification.}\quad
The first and last decimal digits of $p$ are:
\begin{equation}\label{eq:p29998digits}
  p \;=\; 1602087359509060348500037851549319056715\cdots
         16930852697403293697.
\end{equation}

%------------------------------------------------------------------------------
\subsection{Factorisation of $p-1$}
%------------------------------------------------------------------------------

\begin{rigorousbox}
\begin{lemma}[Factorisation of $p-1$]\label{lem:factorisation29998}
With $a, b, m, p$ as in Definition~\ref{def:instance29998},
\[
  p - 1 \;=\; \underbrace{2^{19435}\cdot 3^{19174}}_{F}
              \;\cdot\; m,
\]
where $F := 2^{19435}\cdot 3^{19174}$ has $15\,000$ decimal digits and
satisfies $F > \sqrt{p}$.
\end{lemma}
\end{rigorousbox}

\begin{proof}
Since $m + 1 = 2^{19435}\cdot 3^{19173}$,
\[
  p - 1 \;=\; 3m(m+1) \;=\; 3m\cdot 2^{19435}\cdot 3^{19173}
           \;=\; 2^{19435}\cdot 3^{19174}\cdot m.
\]
For the bound $F > \sqrt{p}$, note that $p < 3(m+1)^2$, so
$\sqrt{p} < \sqrt{3}\,(m+1) < 3(m+1) = 2^{19435}\cdot 3^{19174} = F$.

\noindent Numerically, comparing base-$2$ logarithms:
\[
  2\log_2 F
  \;=\; 2\bigl(19435 + 19174\log_2 3\bigr)
  \;\approx\; 99\,650.14,
\]
\[
  \log_2 p \;\approx\; 99\,648.56,
\]
confirming $F^2 > p$ with a margin of $\approx 1.58$ bits.
\end{proof}

%------------------------------------------------------------------------------
\subsection{The Pocklington--Lehmer Certificate}
%------------------------------------------------------------------------------

Theorem~\ref{thm:Pocklington} (recalled from Section~\ref{sec:family})
requires, for each prime $q \mid F$ — here $q \in \{2, 3\}$ — an integer
$w_q$ satisfying
\[
  w_q^{p-1} \equiv 1 \pmod{p}
  \qquad\text{and}\qquad
  \gcd\!\bigl(w_q^{(p-1)/q} - 1,\; p\bigr) = 1.
\]

\begin{resultbox}[title={\textbf{Certificate} — 29\,998-digit prime}]
\[
  p \;=\; 3m(m+1)+1,\qquad m = 2^{19435}\cdot 3^{19173}-1.
\]

\medskip
\noindent\textbf{Factored part of $p-1$:}
$\;F = 2^{19435}\cdot 3^{19174},\quad F > \sqrt{p}.\quad$\checkmark

\medskip
\renewcommand{\arraystretch}{1.4}
\begin{tabular}{@{}c l l@{}}
\toprule
Prime $q$ & Witness $w_q$ & Conditions verified \\
\midrule
$2$ & $w_2 = 5$ &
  $5^{p-1}\equiv 1\pmod{p}$\quad and\quad
  $\gcd(5^{(p-1)/2}-1,\,p)=1$ \\
$3$ & $w_3 = 7$ &
  $7^{p-1}\equiv 1\pmod{p}$\quad and\quad
  $\gcd(7^{(p-1)/3}-1,\,p)=1$ \\
\bottomrule
\end{tabular}

\medskip
\noindent\textbf{Conclusion:} $p$ is prime (deterministic proof). \checkmark
\end{resultbox}

%------------------------------------------------------------------------------
\subsection{Independent Numerical Verification}
%------------------------------------------------------------------------------

All four conditions of the certificate were verified independently using
\texttt{Python 3.12} with its built-in arbitrary-precision integer
arithmetic (GMP backend). The script performs no probabilistic step;
every operation is exact modular arithmetic.

\begin{rigorousbox}[title={\textbf{Verification record} — independent computation}]
\textbf{Environment:} Python 3.12, \texttt{sys.set\_int\_max\_str\_digits(40000)};
modular exponentiations via \texttt{pow(base, exp, mod)}.\\[4pt]

\textbf{Step 1 — Reconstructing $p$.}\quad
Computed $m = 2^{19435}\cdot 3^{19173}-1$ (14\,999 digits) and
$p = 3m(m+1)+1$ (29\,998 digits) from scratch.
First 40 digits of $p$: \texttt{1602087359509060348500037851549319056715};
last 20 digits: \texttt{16930852697403293697}.
Both match the stored certificate exactly.\\[4pt]

\textbf{Step 2 — Pocklington bound.}
\[
  2\log_2 F = 99\,650.1420 \;>\; \log_2 p = 99\,648.5570. \checkmark
\]

\textbf{Step 3 — Witness $w_2 = 5$ for $q = 2$.}
\begin{align*}
  5^{p-1} &\equiv 1 \pmod{p}, \checkmark \\
  \gcd(5^{(p-1)/2}-1,\,p) &= 1. \checkmark
\end{align*}

\textbf{Step 4 — Witness $w_3 = 7$ for $q = 3$.}
\begin{align*}
  7^{p-1} &\equiv 1 \pmod{p}, \checkmark \\
  \gcd(7^{(p-1)/3}-1,\,p) &= 1. \checkmark
\end{align*}

\textbf{Wall-clock time:} $6\,945$ s $\approx 1$ h $56$ min
(single core, Python GMP, no multi-threading).\\[4pt]

\textbf{Conclusion:} all four Pocklington–Lehmer conditions are satisfied.
The number $p$ is \emph{unconditionally prime}.
\end{rigorousbox}

%------------------------------------------------------------------------------
\subsection{Comparison with the Original Computation}
%------------------------------------------------------------------------------

\begin{table}[h]
\centering
\caption{Original search vs.\ independent certificate verification
         for the 29\,998-digit prime.}
\label{tab:comparison29998}
\renewcommand{\arraystretch}{1.3}
\begin{tabular}{@{}lll@{}}
\toprule
 & \textbf{PrimeQuest v4 (search)} & \textbf{Independent verification} \\
\midrule
Purpose        & Find a prime candidate          & Check the certificate \\
Machine        & ARM64, 9 cores (Snapdragon X)   & Single core (Python GMP) \\
Wall-clock     & $11\,146$ s ($3$ h $06$ min)    & $6\,945$ s ($1$ h $56$ min) \\
MR tests run   & $286$ (probabilistic pre-filter) & — (no probabilistic step) \\
Result         & Certificate produced             & Certificate confirmed \\
\bottomrule
\end{tabular}
\end{table}

The verification is $\approx 1.6\times$ faster than the original search
because it executes exactly four modular exponentiations on the known
candidate, bypassing the sieve, Theorem-2 filter, and Miller–Rabin
pre-screening required during the search phase.

%------------------------------------------------------------------------------
\subsection{Record in Context}
%------------------------------------------------------------------------------

Table~\ref{tab:allcerts} (Section~\ref{sec:results}) lists all four
\textsc{PrimeQuest} certificates. The 29\,998-digit instance is currently
the largest unconditional Pocklington–Lehmer certificate in the family
$p = 3m(m+1)+1$, $m = 2^a\cdot 3^b - 1$, produced on a single consumer
laptop without distributed computing or specialised hardware.

%==============================================================================
% MINIMAL STAND-ALONE PREAMBLE (uncomment for self-contained compilation)
%==============================================================================
% \iffalse   % <-- remove this line to enable stand-alone mode
% \documentclass[11pt,a4paper]{article}
% \usepackage[utf8]{inputenc}
% \usepackage[T1]{fontenc}
% \usepackage[english]{babel}
% \usepackage{amsmath,amssymb,amsthm,mathtools,booktabs}
% \usepackage{geometry}\geometry{margin=2.5cm}
% \usepackage[dvipsnames]{xcolor}
% \usepackage{tcolorbox}
% \tcbuselibrary{skins,breakable}
% \newtcolorbox{rigorousbox}[1][]{colback=Green!8,colframe=Green!60!black,
%   arc=1mm,title=\textbf{Rigorous},fonttitle=\bfseries,breakable,#1}
% \newtcolorbox{resultbox}[1][]{colback=RoyalBlue!5,colframe=RoyalBlue!70!black,
%   arc=1mm,title=\textbf{Computational Result},fonttitle=\bfseries,breakable,#1}
% \newtheorem{theorem}{Theorem}[section]
% \newtheorem{lemma}[theorem]{Lemma}
% \theoremstyle{definition}
% \newtheorem{definition}[theorem]{Definition}
% \begin{document}
% \input{certificate_29998}
% \end{document}
% \fi

% \end{document}
% \fi
          % inside the main document
%
% Self-contained compilation:
%   pdflatex certificate_29998.tex
%
%==============================================================================

%------------------------------------------------------------------------------
\section{Unconditional Certificate for the 29\,998-Digit Prime}
\label{sec:cert29998}
%------------------------------------------------------------------------------

This section provides the complete, independently verified Pocklington–Lehmer
primality certificate for the largest prime currently produced by
\textsc{PrimeQuest}.

%------------------------------------------------------------------------------
\subsection{Parameters and Structure}
%------------------------------------------------------------------------------

\begin{definition}[29\,998-digit PrimeQuest instance]\label{def:instance29998}
Set
\[
  a = 19\,435,\qquad b = 19\,173,
\]
and define
\[
  m \;=\; 2^{19435}\cdot 3^{19173} - 1
    \qquad\bigl(14\,999 \text{ decimal digits}\bigr),
\]
\[
  p \;=\; 3m(m+1)+1
    \qquad\bigl(29\,998 \text{ decimal digits}\bigr).
\]
\end{definition}

The ratio $b/a = 19\,173/19\,435 \approx 0.9865$ is close to the
theoretical optimum
\[
  \rho^* = \frac{\log_{10} 2}{\log_{10} 3} \approx 0.6309
\]
in the sense that it minimises the digit count for a prescribed $a$;
the value used here targets exactly $29\,998$ digits and was located by
the \textsc{PrimeQuest v4} algorithm (see Section~\ref{sec:algo}).

\medskip

\noindent\textbf{Identification.}\quad
The first and last decimal digits of $p$ are:
\begin{equation}\label{eq:p29998digits}
  p \;=\; 1602087359509060348500037851549319056715\cdots
         16930852697403293697.
\end{equation}

%------------------------------------------------------------------------------
\subsection{Factorisation of $p-1$}
%------------------------------------------------------------------------------

\begin{rigorousbox}
\begin{lemma}[Factorisation of $p-1$]\label{lem:factorisation29998}
With $a, b, m, p$ as in Definition~\ref{def:instance29998},
\[
  p - 1 \;=\; \underbrace{2^{19435}\cdot 3^{19174}}_{F}
              \;\cdot\; m,
\]
where $F := 2^{19435}\cdot 3^{19174}$ has $15\,000$ decimal digits and
satisfies $F > \sqrt{p}$.
\end{lemma}
\end{rigorousbox}

\begin{proof}
Since $m + 1 = 2^{19435}\cdot 3^{19173}$,
\[
  p - 1 \;=\; 3m(m+1) \;=\; 3m\cdot 2^{19435}\cdot 3^{19173}
           \;=\; 2^{19435}\cdot 3^{19174}\cdot m.
\]
For the bound $F > \sqrt{p}$, note that $p < 3(m+1)^2$, so
$\sqrt{p} < \sqrt{3}\,(m+1) < 3(m+1) = 2^{19435}\cdot 3^{19174} = F$.

\noindent Numerically, comparing base-$2$ logarithms:
\[
  2\log_2 F
  \;=\; 2\bigl(19435 + 19174\log_2 3\bigr)
  \;\approx\; 99\,650.14,
\]
\[
  \log_2 p \;\approx\; 99\,648.56,
\]
confirming $F^2 > p$ with a margin of $\approx 1.58$ bits.
\end{proof}

%------------------------------------------------------------------------------
\subsection{The Pocklington--Lehmer Certificate}
%------------------------------------------------------------------------------

Theorem~\ref{thm:Pocklington} (recalled from Section~\ref{sec:family})
requires, for each prime $q \mid F$ — here $q \in \{2, 3\}$ — an integer
$w_q$ satisfying
\[
  w_q^{p-1} \equiv 1 \pmod{p}
  \qquad\text{and}\qquad
  \gcd\!\bigl(w_q^{(p-1)/q} - 1,\; p\bigr) = 1.
\]

\begin{resultbox}[title={\textbf{Certificate} — 29\,998-digit prime}]
\[
  p \;=\; 3m(m+1)+1,\qquad m = 2^{19435}\cdot 3^{19173}-1.
\]

\medskip
\noindent\textbf{Factored part of $p-1$:}
$\;F = 2^{19435}\cdot 3^{19174},\quad F > \sqrt{p}.\quad$\checkmark

\medskip
\renewcommand{\arraystretch}{1.4}
\begin{tabular}{@{}c l l@{}}
\toprule
Prime $q$ & Witness $w_q$ & Conditions verified \\
\midrule
$2$ & $w_2 = 5$ &
  $5^{p-1}\equiv 1\pmod{p}$\quad and\quad
  $\gcd(5^{(p-1)/2}-1,\,p)=1$ \\
$3$ & $w_3 = 7$ &
  $7^{p-1}\equiv 1\pmod{p}$\quad and\quad
  $\gcd(7^{(p-1)/3}-1,\,p)=1$ \\
\bottomrule
\end{tabular}

\medskip
\noindent\textbf{Conclusion:} $p$ is prime (deterministic proof). \checkmark
\end{resultbox}

%------------------------------------------------------------------------------
\subsection{Independent Numerical Verification}
%------------------------------------------------------------------------------

All four conditions of the certificate were verified independently using
\texttt{Python 3.12} with its built-in arbitrary-precision integer
arithmetic (GMP backend). The script performs no probabilistic step;
every operation is exact modular arithmetic.

\begin{rigorousbox}[title={\textbf{Verification record} — independent computation}]
\textbf{Environment:} Python 3.12, \texttt{sys.set\_int\_max\_str\_digits(40000)};
modular exponentiations via \texttt{pow(base, exp, mod)}.\\[4pt]

\textbf{Step 1 — Reconstructing $p$.}\quad
Computed $m = 2^{19435}\cdot 3^{19173}-1$ (14\,999 digits) and
$p = 3m(m+1)+1$ (29\,998 digits) from scratch.
First 40 digits of $p$: \texttt{1602087359509060348500037851549319056715};
last 20 digits: \texttt{16930852697403293697}.
Both match the stored certificate exactly.\\[4pt]

\textbf{Step 2 — Pocklington bound.}
\[
  2\log_2 F = 99\,650.1420 \;>\; \log_2 p = 99\,648.5570. \checkmark
\]

\textbf{Step 3 — Witness $w_2 = 5$ for $q = 2$.}
\begin{align*}
  5^{p-1} &\equiv 1 \pmod{p}, \checkmark \\
  \gcd(5^{(p-1)/2}-1,\,p) &= 1. \checkmark
\end{align*}

\textbf{Step 4 — Witness $w_3 = 7$ for $q = 3$.}
\begin{align*}
  7^{p-1} &\equiv 1 \pmod{p}, \checkmark \\
  \gcd(7^{(p-1)/3}-1,\,p) &= 1. \checkmark
\end{align*}

\textbf{Wall-clock time:} $6\,945$ s $\approx 1$ h $56$ min
(single core, Python GMP, no multi-threading).\\[4pt]

\textbf{Conclusion:} all four Pocklington–Lehmer conditions are satisfied.
The number $p$ is \emph{unconditionally prime}.
\end{rigorousbox}

%------------------------------------------------------------------------------
\subsection{Comparison with the Original Computation}
%------------------------------------------------------------------------------

\begin{table}[h]
\centering
\caption{Original search vs.\ independent certificate verification
         for the 29\,998-digit prime.}
\label{tab:comparison29998}
\renewcommand{\arraystretch}{1.3}
\begin{tabular}{@{}lll@{}}
\toprule
 & \textbf{PrimeQuest v4 (search)} & \textbf{Independent verification} \\
\midrule
Purpose        & Find a prime candidate          & Check the certificate \\
Machine        & ARM64, 9 cores (Snapdragon X)   & Single core (Python GMP) \\
Wall-clock     & $11\,146$ s ($3$ h $06$ min)    & $6\,945$ s ($1$ h $56$ min) \\
MR tests run   & $286$ (probabilistic pre-filter) & — (no probabilistic step) \\
Result         & Certificate produced             & Certificate confirmed \\
\bottomrule
\end{tabular}
\end{table}

The verification is $\approx 1.6\times$ faster than the original search
because it executes exactly four modular exponentiations on the known
candidate, bypassing the sieve, Theorem-2 filter, and Miller–Rabin
pre-screening required during the search phase.

%------------------------------------------------------------------------------
\subsection{Record in Context}
%------------------------------------------------------------------------------

Table~\ref{tab:allcerts} (Section~\ref{sec:results}) lists all four
\textsc{PrimeQuest} certificates. The 29\,998-digit instance is currently
the largest unconditional Pocklington–Lehmer certificate in the family
$p = 3m(m+1)+1$, $m = 2^a\cdot 3^b - 1$, produced on a single consumer
laptop without distributed computing or specialised hardware.

%==============================================================================
% MINIMAL STAND-ALONE PREAMBLE (uncomment for self-contained compilation)
%==============================================================================
% \iffalse   % <-- remove this line to enable stand-alone mode
% \documentclass[11pt,a4paper]{article}
% \usepackage[utf8]{inputenc}
% \usepackage[T1]{fontenc}
% \usepackage[english]{babel}
% \usepackage{amsmath,amssymb,amsthm,mathtools,booktabs}
% \usepackage{geometry}\geometry{margin=2.5cm}
% \usepackage[dvipsnames]{xcolor}
% \usepackage{tcolorbox}
% \tcbuselibrary{skins,breakable}
% \newtcolorbox{rigorousbox}[1][]{colback=Green!8,colframe=Green!60!black,
%   arc=1mm,title=\textbf{Rigorous},fonttitle=\bfseries,breakable,#1}
% \newtcolorbox{resultbox}[1][]{colback=RoyalBlue!5,colframe=RoyalBlue!70!black,
%   arc=1mm,title=\textbf{Computational Result},fonttitle=\bfseries,breakable,#1}
% \newtheorem{theorem}{Theorem}[section]
% \newtheorem{lemma}[theorem]{Lemma}
% \theoremstyle{definition}
% \newtheorem{definition}[theorem]{Definition}
% \begin{document}
% %==============================================================================
% Unconditional Primality Certificate — 29 998-digit prime
% Parametric family  p = 3m(m+1)+1,  m = 2^a·3^b − 1
% Independent numerical verification — May 2026
%==============================================================================
%
% Stand-alone usage (requires paper_pocklington.tex preamble, or use the
% minimal preamble at the bottom of this file):
%
%   %==============================================================================
% Unconditional Primality Certificate — 29 998-digit prime
% Parametric family  p = 3m(m+1)+1,  m = 2^a·3^b − 1
% Independent numerical verification — May 2026
%==============================================================================
%
% Stand-alone usage (requires paper_pocklington.tex preamble, or use the
% minimal preamble at the bottom of this file):
%
%   \input{certificate_29998}          % inside the main document
%
% Self-contained compilation:
%   pdflatex certificate_29998.tex
%
%==============================================================================

%------------------------------------------------------------------------------
\section{Unconditional Certificate for the 29\,998-Digit Prime}
\label{sec:cert29998}
%------------------------------------------------------------------------------

This section provides the complete, independently verified Pocklington–Lehmer
primality certificate for the largest prime currently produced by
\textsc{PrimeQuest}.

%------------------------------------------------------------------------------
\subsection{Parameters and Structure}
%------------------------------------------------------------------------------

\begin{definition}[29\,998-digit PrimeQuest instance]\label{def:instance29998}
Set
\[
  a = 19\,435,\qquad b = 19\,173,
\]
and define
\[
  m \;=\; 2^{19435}\cdot 3^{19173} - 1
    \qquad\bigl(14\,999 \text{ decimal digits}\bigr),
\]
\[
  p \;=\; 3m(m+1)+1
    \qquad\bigl(29\,998 \text{ decimal digits}\bigr).
\]
\end{definition}

The ratio $b/a = 19\,173/19\,435 \approx 0.9865$ is close to the
theoretical optimum
\[
  \rho^* = \frac{\log_{10} 2}{\log_{10} 3} \approx 0.6309
\]
in the sense that it minimises the digit count for a prescribed $a$;
the value used here targets exactly $29\,998$ digits and was located by
the \textsc{PrimeQuest v4} algorithm (see Section~\ref{sec:algo}).

\medskip

\noindent\textbf{Identification.}\quad
The first and last decimal digits of $p$ are:
\begin{equation}\label{eq:p29998digits}
  p \;=\; 1602087359509060348500037851549319056715\cdots
         16930852697403293697.
\end{equation}

%------------------------------------------------------------------------------
\subsection{Factorisation of $p-1$}
%------------------------------------------------------------------------------

\begin{rigorousbox}
\begin{lemma}[Factorisation of $p-1$]\label{lem:factorisation29998}
With $a, b, m, p$ as in Definition~\ref{def:instance29998},
\[
  p - 1 \;=\; \underbrace{2^{19435}\cdot 3^{19174}}_{F}
              \;\cdot\; m,
\]
where $F := 2^{19435}\cdot 3^{19174}$ has $15\,000$ decimal digits and
satisfies $F > \sqrt{p}$.
\end{lemma}
\end{rigorousbox}

\begin{proof}
Since $m + 1 = 2^{19435}\cdot 3^{19173}$,
\[
  p - 1 \;=\; 3m(m+1) \;=\; 3m\cdot 2^{19435}\cdot 3^{19173}
           \;=\; 2^{19435}\cdot 3^{19174}\cdot m.
\]
For the bound $F > \sqrt{p}$, note that $p < 3(m+1)^2$, so
$\sqrt{p} < \sqrt{3}\,(m+1) < 3(m+1) = 2^{19435}\cdot 3^{19174} = F$.

\noindent Numerically, comparing base-$2$ logarithms:
\[
  2\log_2 F
  \;=\; 2\bigl(19435 + 19174\log_2 3\bigr)
  \;\approx\; 99\,650.14,
\]
\[
  \log_2 p \;\approx\; 99\,648.56,
\]
confirming $F^2 > p$ with a margin of $\approx 1.58$ bits.
\end{proof}

%------------------------------------------------------------------------------
\subsection{The Pocklington--Lehmer Certificate}
%------------------------------------------------------------------------------

Theorem~\ref{thm:Pocklington} (recalled from Section~\ref{sec:family})
requires, for each prime $q \mid F$ — here $q \in \{2, 3\}$ — an integer
$w_q$ satisfying
\[
  w_q^{p-1} \equiv 1 \pmod{p}
  \qquad\text{and}\qquad
  \gcd\!\bigl(w_q^{(p-1)/q} - 1,\; p\bigr) = 1.
\]

\begin{resultbox}[title={\textbf{Certificate} — 29\,998-digit prime}]
\[
  p \;=\; 3m(m+1)+1,\qquad m = 2^{19435}\cdot 3^{19173}-1.
\]

\medskip
\noindent\textbf{Factored part of $p-1$:}
$\;F = 2^{19435}\cdot 3^{19174},\quad F > \sqrt{p}.\quad$\checkmark

\medskip
\renewcommand{\arraystretch}{1.4}
\begin{tabular}{@{}c l l@{}}
\toprule
Prime $q$ & Witness $w_q$ & Conditions verified \\
\midrule
$2$ & $w_2 = 5$ &
  $5^{p-1}\equiv 1\pmod{p}$\quad and\quad
  $\gcd(5^{(p-1)/2}-1,\,p)=1$ \\
$3$ & $w_3 = 7$ &
  $7^{p-1}\equiv 1\pmod{p}$\quad and\quad
  $\gcd(7^{(p-1)/3}-1,\,p)=1$ \\
\bottomrule
\end{tabular}

\medskip
\noindent\textbf{Conclusion:} $p$ is prime (deterministic proof). \checkmark
\end{resultbox}

%------------------------------------------------------------------------------
\subsection{Independent Numerical Verification}
%------------------------------------------------------------------------------

All four conditions of the certificate were verified independently using
\texttt{Python 3.12} with its built-in arbitrary-precision integer
arithmetic (GMP backend). The script performs no probabilistic step;
every operation is exact modular arithmetic.

\begin{rigorousbox}[title={\textbf{Verification record} — independent computation}]
\textbf{Environment:} Python 3.12, \texttt{sys.set\_int\_max\_str\_digits(40000)};
modular exponentiations via \texttt{pow(base, exp, mod)}.\\[4pt]

\textbf{Step 1 — Reconstructing $p$.}\quad
Computed $m = 2^{19435}\cdot 3^{19173}-1$ (14\,999 digits) and
$p = 3m(m+1)+1$ (29\,998 digits) from scratch.
First 40 digits of $p$: \texttt{1602087359509060348500037851549319056715};
last 20 digits: \texttt{16930852697403293697}.
Both match the stored certificate exactly.\\[4pt]

\textbf{Step 2 — Pocklington bound.}
\[
  2\log_2 F = 99\,650.1420 \;>\; \log_2 p = 99\,648.5570. \checkmark
\]

\textbf{Step 3 — Witness $w_2 = 5$ for $q = 2$.}
\begin{align*}
  5^{p-1} &\equiv 1 \pmod{p}, \checkmark \\
  \gcd(5^{(p-1)/2}-1,\,p) &= 1. \checkmark
\end{align*}

\textbf{Step 4 — Witness $w_3 = 7$ for $q = 3$.}
\begin{align*}
  7^{p-1} &\equiv 1 \pmod{p}, \checkmark \\
  \gcd(7^{(p-1)/3}-1,\,p) &= 1. \checkmark
\end{align*}

\textbf{Wall-clock time:} $6\,945$ s $\approx 1$ h $56$ min
(single core, Python GMP, no multi-threading).\\[4pt]

\textbf{Conclusion:} all four Pocklington–Lehmer conditions are satisfied.
The number $p$ is \emph{unconditionally prime}.
\end{rigorousbox}

%------------------------------------------------------------------------------
\subsection{Comparison with the Original Computation}
%------------------------------------------------------------------------------

\begin{table}[h]
\centering
\caption{Original search vs.\ independent certificate verification
         for the 29\,998-digit prime.}
\label{tab:comparison29998}
\renewcommand{\arraystretch}{1.3}
\begin{tabular}{@{}lll@{}}
\toprule
 & \textbf{PrimeQuest v4 (search)} & \textbf{Independent verification} \\
\midrule
Purpose        & Find a prime candidate          & Check the certificate \\
Machine        & ARM64, 9 cores (Snapdragon X)   & Single core (Python GMP) \\
Wall-clock     & $11\,146$ s ($3$ h $06$ min)    & $6\,945$ s ($1$ h $56$ min) \\
MR tests run   & $286$ (probabilistic pre-filter) & — (no probabilistic step) \\
Result         & Certificate produced             & Certificate confirmed \\
\bottomrule
\end{tabular}
\end{table}

The verification is $\approx 1.6\times$ faster than the original search
because it executes exactly four modular exponentiations on the known
candidate, bypassing the sieve, Theorem-2 filter, and Miller–Rabin
pre-screening required during the search phase.

%------------------------------------------------------------------------------
\subsection{Record in Context}
%------------------------------------------------------------------------------

Table~\ref{tab:allcerts} (Section~\ref{sec:results}) lists all four
\textsc{PrimeQuest} certificates. The 29\,998-digit instance is currently
the largest unconditional Pocklington–Lehmer certificate in the family
$p = 3m(m+1)+1$, $m = 2^a\cdot 3^b - 1$, produced on a single consumer
laptop without distributed computing or specialised hardware.

%==============================================================================
% MINIMAL STAND-ALONE PREAMBLE (uncomment for self-contained compilation)
%==============================================================================
% \iffalse   % <-- remove this line to enable stand-alone mode
% \documentclass[11pt,a4paper]{article}
% \usepackage[utf8]{inputenc}
% \usepackage[T1]{fontenc}
% \usepackage[english]{babel}
% \usepackage{amsmath,amssymb,amsthm,mathtools,booktabs}
% \usepackage{geometry}\geometry{margin=2.5cm}
% \usepackage[dvipsnames]{xcolor}
% \usepackage{tcolorbox}
% \tcbuselibrary{skins,breakable}
% \newtcolorbox{rigorousbox}[1][]{colback=Green!8,colframe=Green!60!black,
%   arc=1mm,title=\textbf{Rigorous},fonttitle=\bfseries,breakable,#1}
% \newtcolorbox{resultbox}[1][]{colback=RoyalBlue!5,colframe=RoyalBlue!70!black,
%   arc=1mm,title=\textbf{Computational Result},fonttitle=\bfseries,breakable,#1}
% \newtheorem{theorem}{Theorem}[section]
% \newtheorem{lemma}[theorem]{Lemma}
% \theoremstyle{definition}
% \newtheorem{definition}[theorem]{Definition}
% \begin{document}
% \input{certificate_29998}
% \end{document}
% \fi
          % inside the main document
%
% Self-contained compilation:
%   pdflatex certificate_29998.tex
%
%==============================================================================

%------------------------------------------------------------------------------
\section{Unconditional Certificate for the 29\,998-Digit Prime}
\label{sec:cert29998}
%------------------------------------------------------------------------------

This section provides the complete, independently verified Pocklington–Lehmer
primality certificate for the largest prime currently produced by
\textsc{PrimeQuest}.

%------------------------------------------------------------------------------
\subsection{Parameters and Structure}
%------------------------------------------------------------------------------

\begin{definition}[29\,998-digit PrimeQuest instance]\label{def:instance29998}
Set
\[
  a = 19\,435,\qquad b = 19\,173,
\]
and define
\[
  m \;=\; 2^{19435}\cdot 3^{19173} - 1
    \qquad\bigl(14\,999 \text{ decimal digits}\bigr),
\]
\[
  p \;=\; 3m(m+1)+1
    \qquad\bigl(29\,998 \text{ decimal digits}\bigr).
\]
\end{definition}

The ratio $b/a = 19\,173/19\,435 \approx 0.9865$ is close to the
theoretical optimum
\[
  \rho^* = \frac{\log_{10} 2}{\log_{10} 3} \approx 0.6309
\]
in the sense that it minimises the digit count for a prescribed $a$;
the value used here targets exactly $29\,998$ digits and was located by
the \textsc{PrimeQuest v4} algorithm (see Section~\ref{sec:algo}).

\medskip

\noindent\textbf{Identification.}\quad
The first and last decimal digits of $p$ are:
\begin{equation}\label{eq:p29998digits}
  p \;=\; 1602087359509060348500037851549319056715\cdots
         16930852697403293697.
\end{equation}

%------------------------------------------------------------------------------
\subsection{Factorisation of $p-1$}
%------------------------------------------------------------------------------

\begin{rigorousbox}
\begin{lemma}[Factorisation of $p-1$]\label{lem:factorisation29998}
With $a, b, m, p$ as in Definition~\ref{def:instance29998},
\[
  p - 1 \;=\; \underbrace{2^{19435}\cdot 3^{19174}}_{F}
              \;\cdot\; m,
\]
where $F := 2^{19435}\cdot 3^{19174}$ has $15\,000$ decimal digits and
satisfies $F > \sqrt{p}$.
\end{lemma}
\end{rigorousbox}

\begin{proof}
Since $m + 1 = 2^{19435}\cdot 3^{19173}$,
\[
  p - 1 \;=\; 3m(m+1) \;=\; 3m\cdot 2^{19435}\cdot 3^{19173}
           \;=\; 2^{19435}\cdot 3^{19174}\cdot m.
\]
For the bound $F > \sqrt{p}$, note that $p < 3(m+1)^2$, so
$\sqrt{p} < \sqrt{3}\,(m+1) < 3(m+1) = 2^{19435}\cdot 3^{19174} = F$.

\noindent Numerically, comparing base-$2$ logarithms:
\[
  2\log_2 F
  \;=\; 2\bigl(19435 + 19174\log_2 3\bigr)
  \;\approx\; 99\,650.14,
\]
\[
  \log_2 p \;\approx\; 99\,648.56,
\]
confirming $F^2 > p$ with a margin of $\approx 1.58$ bits.
\end{proof}

%------------------------------------------------------------------------------
\subsection{The Pocklington--Lehmer Certificate}
%------------------------------------------------------------------------------

Theorem~\ref{thm:Pocklington} (recalled from Section~\ref{sec:family})
requires, for each prime $q \mid F$ — here $q \in \{2, 3\}$ — an integer
$w_q$ satisfying
\[
  w_q^{p-1} \equiv 1 \pmod{p}
  \qquad\text{and}\qquad
  \gcd\!\bigl(w_q^{(p-1)/q} - 1,\; p\bigr) = 1.
\]

\begin{resultbox}[title={\textbf{Certificate} — 29\,998-digit prime}]
\[
  p \;=\; 3m(m+1)+1,\qquad m = 2^{19435}\cdot 3^{19173}-1.
\]

\medskip
\noindent\textbf{Factored part of $p-1$:}
$\;F = 2^{19435}\cdot 3^{19174},\quad F > \sqrt{p}.\quad$\checkmark

\medskip
\renewcommand{\arraystretch}{1.4}
\begin{tabular}{@{}c l l@{}}
\toprule
Prime $q$ & Witness $w_q$ & Conditions verified \\
\midrule
$2$ & $w_2 = 5$ &
  $5^{p-1}\equiv 1\pmod{p}$\quad and\quad
  $\gcd(5^{(p-1)/2}-1,\,p)=1$ \\
$3$ & $w_3 = 7$ &
  $7^{p-1}\equiv 1\pmod{p}$\quad and\quad
  $\gcd(7^{(p-1)/3}-1,\,p)=1$ \\
\bottomrule
\end{tabular}

\medskip
\noindent\textbf{Conclusion:} $p$ is prime (deterministic proof). \checkmark
\end{resultbox}

%------------------------------------------------------------------------------
\subsection{Independent Numerical Verification}
%------------------------------------------------------------------------------

All four conditions of the certificate were verified independently using
\texttt{Python 3.12} with its built-in arbitrary-precision integer
arithmetic (GMP backend). The script performs no probabilistic step;
every operation is exact modular arithmetic.

\begin{rigorousbox}[title={\textbf{Verification record} — independent computation}]
\textbf{Environment:} Python 3.12, \texttt{sys.set\_int\_max\_str\_digits(40000)};
modular exponentiations via \texttt{pow(base, exp, mod)}.\\[4pt]

\textbf{Step 1 — Reconstructing $p$.}\quad
Computed $m = 2^{19435}\cdot 3^{19173}-1$ (14\,999 digits) and
$p = 3m(m+1)+1$ (29\,998 digits) from scratch.
First 40 digits of $p$: \texttt{1602087359509060348500037851549319056715};
last 20 digits: \texttt{16930852697403293697}.
Both match the stored certificate exactly.\\[4pt]

\textbf{Step 2 — Pocklington bound.}
\[
  2\log_2 F = 99\,650.1420 \;>\; \log_2 p = 99\,648.5570. \checkmark
\]

\textbf{Step 3 — Witness $w_2 = 5$ for $q = 2$.}
\begin{align*}
  5^{p-1} &\equiv 1 \pmod{p}, \checkmark \\
  \gcd(5^{(p-1)/2}-1,\,p) &= 1. \checkmark
\end{align*}

\textbf{Step 4 — Witness $w_3 = 7$ for $q = 3$.}
\begin{align*}
  7^{p-1} &\equiv 1 \pmod{p}, \checkmark \\
  \gcd(7^{(p-1)/3}-1,\,p) &= 1. \checkmark
\end{align*}

\textbf{Wall-clock time:} $6\,945$ s $\approx 1$ h $56$ min
(single core, Python GMP, no multi-threading).\\[4pt]

\textbf{Conclusion:} all four Pocklington–Lehmer conditions are satisfied.
The number $p$ is \emph{unconditionally prime}.
\end{rigorousbox}

%------------------------------------------------------------------------------
\subsection{Comparison with the Original Computation}
%------------------------------------------------------------------------------

\begin{table}[h]
\centering
\caption{Original search vs.\ independent certificate verification
         for the 29\,998-digit prime.}
\label{tab:comparison29998}
\renewcommand{\arraystretch}{1.3}
\begin{tabular}{@{}lll@{}}
\toprule
 & \textbf{PrimeQuest v4 (search)} & \textbf{Independent verification} \\
\midrule
Purpose        & Find a prime candidate          & Check the certificate \\
Machine        & ARM64, 9 cores (Snapdragon X)   & Single core (Python GMP) \\
Wall-clock     & $11\,146$ s ($3$ h $06$ min)    & $6\,945$ s ($1$ h $56$ min) \\
MR tests run   & $286$ (probabilistic pre-filter) & — (no probabilistic step) \\
Result         & Certificate produced             & Certificate confirmed \\
\bottomrule
\end{tabular}
\end{table}

The verification is $\approx 1.6\times$ faster than the original search
because it executes exactly four modular exponentiations on the known
candidate, bypassing the sieve, Theorem-2 filter, and Miller–Rabin
pre-screening required during the search phase.

%------------------------------------------------------------------------------
\subsection{Record in Context}
%------------------------------------------------------------------------------

Table~\ref{tab:allcerts} (Section~\ref{sec:results}) lists all four
\textsc{PrimeQuest} certificates. The 29\,998-digit instance is currently
the largest unconditional Pocklington–Lehmer certificate in the family
$p = 3m(m+1)+1$, $m = 2^a\cdot 3^b - 1$, produced on a single consumer
laptop without distributed computing or specialised hardware.

%==============================================================================
% MINIMAL STAND-ALONE PREAMBLE (uncomment for self-contained compilation)
%==============================================================================
% \iffalse   % <-- remove this line to enable stand-alone mode
% \documentclass[11pt,a4paper]{article}
% \usepackage[utf8]{inputenc}
% \usepackage[T1]{fontenc}
% \usepackage[english]{babel}
% \usepackage{amsmath,amssymb,amsthm,mathtools,booktabs}
% \usepackage{geometry}\geometry{margin=2.5cm}
% \usepackage[dvipsnames]{xcolor}
% \usepackage{tcolorbox}
% \tcbuselibrary{skins,breakable}
% \newtcolorbox{rigorousbox}[1][]{colback=Green!8,colframe=Green!60!black,
%   arc=1mm,title=\textbf{Rigorous},fonttitle=\bfseries,breakable,#1}
% \newtcolorbox{resultbox}[1][]{colback=RoyalBlue!5,colframe=RoyalBlue!70!black,
%   arc=1mm,title=\textbf{Computational Result},fonttitle=\bfseries,breakable,#1}
% \newtheorem{theorem}{Theorem}[section]
% \newtheorem{lemma}[theorem]{Lemma}
% \theoremstyle{definition}
% \newtheorem{definition}[theorem]{Definition}
% \begin{document}
% %==============================================================================
% Unconditional Primality Certificate — 29 998-digit prime
% Parametric family  p = 3m(m+1)+1,  m = 2^a·3^b − 1
% Independent numerical verification — May 2026
%==============================================================================
%
% Stand-alone usage (requires paper_pocklington.tex preamble, or use the
% minimal preamble at the bottom of this file):
%
%   \input{certificate_29998}          % inside the main document
%
% Self-contained compilation:
%   pdflatex certificate_29998.tex
%
%==============================================================================

%------------------------------------------------------------------------------
\section{Unconditional Certificate for the 29\,998-Digit Prime}
\label{sec:cert29998}
%------------------------------------------------------------------------------

This section provides the complete, independently verified Pocklington–Lehmer
primality certificate for the largest prime currently produced by
\textsc{PrimeQuest}.

%------------------------------------------------------------------------------
\subsection{Parameters and Structure}
%------------------------------------------------------------------------------

\begin{definition}[29\,998-digit PrimeQuest instance]\label{def:instance29998}
Set
\[
  a = 19\,435,\qquad b = 19\,173,
\]
and define
\[
  m \;=\; 2^{19435}\cdot 3^{19173} - 1
    \qquad\bigl(14\,999 \text{ decimal digits}\bigr),
\]
\[
  p \;=\; 3m(m+1)+1
    \qquad\bigl(29\,998 \text{ decimal digits}\bigr).
\]
\end{definition}

The ratio $b/a = 19\,173/19\,435 \approx 0.9865$ is close to the
theoretical optimum
\[
  \rho^* = \frac{\log_{10} 2}{\log_{10} 3} \approx 0.6309
\]
in the sense that it minimises the digit count for a prescribed $a$;
the value used here targets exactly $29\,998$ digits and was located by
the \textsc{PrimeQuest v4} algorithm (see Section~\ref{sec:algo}).

\medskip

\noindent\textbf{Identification.}\quad
The first and last decimal digits of $p$ are:
\begin{equation}\label{eq:p29998digits}
  p \;=\; 1602087359509060348500037851549319056715\cdots
         16930852697403293697.
\end{equation}

%------------------------------------------------------------------------------
\subsection{Factorisation of $p-1$}
%------------------------------------------------------------------------------

\begin{rigorousbox}
\begin{lemma}[Factorisation of $p-1$]\label{lem:factorisation29998}
With $a, b, m, p$ as in Definition~\ref{def:instance29998},
\[
  p - 1 \;=\; \underbrace{2^{19435}\cdot 3^{19174}}_{F}
              \;\cdot\; m,
\]
where $F := 2^{19435}\cdot 3^{19174}$ has $15\,000$ decimal digits and
satisfies $F > \sqrt{p}$.
\end{lemma}
\end{rigorousbox}

\begin{proof}
Since $m + 1 = 2^{19435}\cdot 3^{19173}$,
\[
  p - 1 \;=\; 3m(m+1) \;=\; 3m\cdot 2^{19435}\cdot 3^{19173}
           \;=\; 2^{19435}\cdot 3^{19174}\cdot m.
\]
For the bound $F > \sqrt{p}$, note that $p < 3(m+1)^2$, so
$\sqrt{p} < \sqrt{3}\,(m+1) < 3(m+1) = 2^{19435}\cdot 3^{19174} = F$.

\noindent Numerically, comparing base-$2$ logarithms:
\[
  2\log_2 F
  \;=\; 2\bigl(19435 + 19174\log_2 3\bigr)
  \;\approx\; 99\,650.14,
\]
\[
  \log_2 p \;\approx\; 99\,648.56,
\]
confirming $F^2 > p$ with a margin of $\approx 1.58$ bits.
\end{proof}

%------------------------------------------------------------------------------
\subsection{The Pocklington--Lehmer Certificate}
%------------------------------------------------------------------------------

Theorem~\ref{thm:Pocklington} (recalled from Section~\ref{sec:family})
requires, for each prime $q \mid F$ — here $q \in \{2, 3\}$ — an integer
$w_q$ satisfying
\[
  w_q^{p-1} \equiv 1 \pmod{p}
  \qquad\text{and}\qquad
  \gcd\!\bigl(w_q^{(p-1)/q} - 1,\; p\bigr) = 1.
\]

\begin{resultbox}[title={\textbf{Certificate} — 29\,998-digit prime}]
\[
  p \;=\; 3m(m+1)+1,\qquad m = 2^{19435}\cdot 3^{19173}-1.
\]

\medskip
\noindent\textbf{Factored part of $p-1$:}
$\;F = 2^{19435}\cdot 3^{19174},\quad F > \sqrt{p}.\quad$\checkmark

\medskip
\renewcommand{\arraystretch}{1.4}
\begin{tabular}{@{}c l l@{}}
\toprule
Prime $q$ & Witness $w_q$ & Conditions verified \\
\midrule
$2$ & $w_2 = 5$ &
  $5^{p-1}\equiv 1\pmod{p}$\quad and\quad
  $\gcd(5^{(p-1)/2}-1,\,p)=1$ \\
$3$ & $w_3 = 7$ &
  $7^{p-1}\equiv 1\pmod{p}$\quad and\quad
  $\gcd(7^{(p-1)/3}-1,\,p)=1$ \\
\bottomrule
\end{tabular}

\medskip
\noindent\textbf{Conclusion:} $p$ is prime (deterministic proof). \checkmark
\end{resultbox}

%------------------------------------------------------------------------------
\subsection{Independent Numerical Verification}
%------------------------------------------------------------------------------

All four conditions of the certificate were verified independently using
\texttt{Python 3.12} with its built-in arbitrary-precision integer
arithmetic (GMP backend). The script performs no probabilistic step;
every operation is exact modular arithmetic.

\begin{rigorousbox}[title={\textbf{Verification record} — independent computation}]
\textbf{Environment:} Python 3.12, \texttt{sys.set\_int\_max\_str\_digits(40000)};
modular exponentiations via \texttt{pow(base, exp, mod)}.\\[4pt]

\textbf{Step 1 — Reconstructing $p$.}\quad
Computed $m = 2^{19435}\cdot 3^{19173}-1$ (14\,999 digits) and
$p = 3m(m+1)+1$ (29\,998 digits) from scratch.
First 40 digits of $p$: \texttt{1602087359509060348500037851549319056715};
last 20 digits: \texttt{16930852697403293697}.
Both match the stored certificate exactly.\\[4pt]

\textbf{Step 2 — Pocklington bound.}
\[
  2\log_2 F = 99\,650.1420 \;>\; \log_2 p = 99\,648.5570. \checkmark
\]

\textbf{Step 3 — Witness $w_2 = 5$ for $q = 2$.}
\begin{align*}
  5^{p-1} &\equiv 1 \pmod{p}, \checkmark \\
  \gcd(5^{(p-1)/2}-1,\,p) &= 1. \checkmark
\end{align*}

\textbf{Step 4 — Witness $w_3 = 7$ for $q = 3$.}
\begin{align*}
  7^{p-1} &\equiv 1 \pmod{p}, \checkmark \\
  \gcd(7^{(p-1)/3}-1,\,p) &= 1. \checkmark
\end{align*}

\textbf{Wall-clock time:} $6\,945$ s $\approx 1$ h $56$ min
(single core, Python GMP, no multi-threading).\\[4pt]

\textbf{Conclusion:} all four Pocklington–Lehmer conditions are satisfied.
The number $p$ is \emph{unconditionally prime}.
\end{rigorousbox}

%------------------------------------------------------------------------------
\subsection{Comparison with the Original Computation}
%------------------------------------------------------------------------------

\begin{table}[h]
\centering
\caption{Original search vs.\ independent certificate verification
         for the 29\,998-digit prime.}
\label{tab:comparison29998}
\renewcommand{\arraystretch}{1.3}
\begin{tabular}{@{}lll@{}}
\toprule
 & \textbf{PrimeQuest v4 (search)} & \textbf{Independent verification} \\
\midrule
Purpose        & Find a prime candidate          & Check the certificate \\
Machine        & ARM64, 9 cores (Snapdragon X)   & Single core (Python GMP) \\
Wall-clock     & $11\,146$ s ($3$ h $06$ min)    & $6\,945$ s ($1$ h $56$ min) \\
MR tests run   & $286$ (probabilistic pre-filter) & — (no probabilistic step) \\
Result         & Certificate produced             & Certificate confirmed \\
\bottomrule
\end{tabular}
\end{table}

The verification is $\approx 1.6\times$ faster than the original search
because it executes exactly four modular exponentiations on the known
candidate, bypassing the sieve, Theorem-2 filter, and Miller–Rabin
pre-screening required during the search phase.

%------------------------------------------------------------------------------
\subsection{Record in Context}
%------------------------------------------------------------------------------

Table~\ref{tab:allcerts} (Section~\ref{sec:results}) lists all four
\textsc{PrimeQuest} certificates. The 29\,998-digit instance is currently
the largest unconditional Pocklington–Lehmer certificate in the family
$p = 3m(m+1)+1$, $m = 2^a\cdot 3^b - 1$, produced on a single consumer
laptop without distributed computing or specialised hardware.

%==============================================================================
% MINIMAL STAND-ALONE PREAMBLE (uncomment for self-contained compilation)
%==============================================================================
% \iffalse   % <-- remove this line to enable stand-alone mode
% \documentclass[11pt,a4paper]{article}
% \usepackage[utf8]{inputenc}
% \usepackage[T1]{fontenc}
% \usepackage[english]{babel}
% \usepackage{amsmath,amssymb,amsthm,mathtools,booktabs}
% \usepackage{geometry}\geometry{margin=2.5cm}
% \usepackage[dvipsnames]{xcolor}
% \usepackage{tcolorbox}
% \tcbuselibrary{skins,breakable}
% \newtcolorbox{rigorousbox}[1][]{colback=Green!8,colframe=Green!60!black,
%   arc=1mm,title=\textbf{Rigorous},fonttitle=\bfseries,breakable,#1}
% \newtcolorbox{resultbox}[1][]{colback=RoyalBlue!5,colframe=RoyalBlue!70!black,
%   arc=1mm,title=\textbf{Computational Result},fonttitle=\bfseries,breakable,#1}
% \newtheorem{theorem}{Theorem}[section]
% \newtheorem{lemma}[theorem]{Lemma}
% \theoremstyle{definition}
% \newtheorem{definition}[theorem]{Definition}
% \begin{document}
% \input{certificate_29998}
% \end{document}
% \fi

% \end{document}
% \fi

% \end{document}
% \fi
          % inside the main document
%
% Self-contained compilation:
%   pdflatex certificate_29998.tex
%
%==============================================================================

%------------------------------------------------------------------------------
\section{Unconditional Certificate for the 29\,998-Digit Prime}
\label{sec:cert29998}
%------------------------------------------------------------------------------

This section provides the complete, independently verified Pocklington–Lehmer
primality certificate for the largest prime currently produced by
\textsc{PrimeQuest}.

%------------------------------------------------------------------------------
\subsection{Parameters and Structure}
%------------------------------------------------------------------------------

\begin{definition}[29\,998-digit PrimeQuest instance]\label{def:instance29998}
Set
\[
  a = 19\,435,\qquad b = 19\,173,
\]
and define
\[
  m \;=\; 2^{19435}\cdot 3^{19173} - 1
    \qquad\bigl(14\,999 \text{ decimal digits}\bigr),
\]
\[
  p \;=\; 3m(m+1)+1
    \qquad\bigl(29\,998 \text{ decimal digits}\bigr).
\]
\end{definition}

The ratio $b/a = 19\,173/19\,435 \approx 0.9865$ is close to the
theoretical optimum
\[
  \rho^* = \frac{\log_{10} 2}{\log_{10} 3} \approx 0.6309
\]
in the sense that it minimises the digit count for a prescribed $a$;
the value used here targets exactly $29\,998$ digits and was located by
the \textsc{PrimeQuest v4} algorithm (see Section~\ref{sec:algo}).

\medskip

\noindent\textbf{Identification.}\quad
The first and last decimal digits of $p$ are:
\begin{equation}\label{eq:p29998digits}
  p \;=\; 1602087359509060348500037851549319056715\cdots
         16930852697403293697.
\end{equation}

%------------------------------------------------------------------------------
\subsection{Factorisation of $p-1$}
%------------------------------------------------------------------------------

\begin{rigorousbox}
\begin{lemma}[Factorisation of $p-1$]\label{lem:factorisation29998}
With $a, b, m, p$ as in Definition~\ref{def:instance29998},
\[
  p - 1 \;=\; \underbrace{2^{19435}\cdot 3^{19174}}_{F}
              \;\cdot\; m,
\]
where $F := 2^{19435}\cdot 3^{19174}$ has $15\,000$ decimal digits and
satisfies $F > \sqrt{p}$.
\end{lemma}
\end{rigorousbox}

\begin{proof}
Since $m + 1 = 2^{19435}\cdot 3^{19173}$,
\[
  p - 1 \;=\; 3m(m+1) \;=\; 3m\cdot 2^{19435}\cdot 3^{19173}
           \;=\; 2^{19435}\cdot 3^{19174}\cdot m.
\]
For the bound $F > \sqrt{p}$, note that $p < 3(m+1)^2$, so
$\sqrt{p} < \sqrt{3}\,(m+1) < 3(m+1) = 2^{19435}\cdot 3^{19174} = F$.

\noindent Numerically, comparing base-$2$ logarithms:
\[
  2\log_2 F
  \;=\; 2\bigl(19435 + 19174\log_2 3\bigr)
  \;\approx\; 99\,650.14,
\]
\[
  \log_2 p \;\approx\; 99\,648.56,
\]
confirming $F^2 > p$ with a margin of $\approx 1.58$ bits.
\end{proof}

%------------------------------------------------------------------------------
\subsection{The Pocklington--Lehmer Certificate}
%------------------------------------------------------------------------------

Theorem~\ref{thm:Pocklington} (recalled from Section~\ref{sec:family})
requires, for each prime $q \mid F$ — here $q \in \{2, 3\}$ — an integer
$w_q$ satisfying
\[
  w_q^{p-1} \equiv 1 \pmod{p}
  \qquad\text{and}\qquad
  \gcd\!\bigl(w_q^{(p-1)/q} - 1,\; p\bigr) = 1.
\]

\begin{resultbox}[title={\textbf{Certificate} — 29\,998-digit prime}]
\[
  p \;=\; 3m(m+1)+1,\qquad m = 2^{19435}\cdot 3^{19173}-1.
\]

\medskip
\noindent\textbf{Factored part of $p-1$:}
$\;F = 2^{19435}\cdot 3^{19174},\quad F > \sqrt{p}.\quad$\checkmark

\medskip
\renewcommand{\arraystretch}{1.4}
\begin{tabular}{@{}c l l@{}}
\toprule
Prime $q$ & Witness $w_q$ & Conditions verified \\
\midrule
$2$ & $w_2 = 5$ &
  $5^{p-1}\equiv 1\pmod{p}$\quad and\quad
  $\gcd(5^{(p-1)/2}-1,\,p)=1$ \\
$3$ & $w_3 = 7$ &
  $7^{p-1}\equiv 1\pmod{p}$\quad and\quad
  $\gcd(7^{(p-1)/3}-1,\,p)=1$ \\
\bottomrule
\end{tabular}

\medskip
\noindent\textbf{Conclusion:} $p$ is prime (deterministic proof). \checkmark
\end{resultbox}

%------------------------------------------------------------------------------
\subsection{Independent Numerical Verification}
%------------------------------------------------------------------------------

All four conditions of the certificate were verified independently using
\texttt{Python 3.12} with its built-in arbitrary-precision integer
arithmetic (GMP backend). The script performs no probabilistic step;
every operation is exact modular arithmetic.

\begin{rigorousbox}[title={\textbf{Verification record} — independent computation}]
\textbf{Environment:} Python 3.12, \texttt{sys.set\_int\_max\_str\_digits(40000)};
modular exponentiations via \texttt{pow(base, exp, mod)}.\\[4pt]

\textbf{Step 1 — Reconstructing $p$.}\quad
Computed $m = 2^{19435}\cdot 3^{19173}-1$ (14\,999 digits) and
$p = 3m(m+1)+1$ (29\,998 digits) from scratch.
First 40 digits of $p$: \texttt{1602087359509060348500037851549319056715};
last 20 digits: \texttt{16930852697403293697}.
Both match the stored certificate exactly.\\[4pt]

\textbf{Step 2 — Pocklington bound.}
\[
  2\log_2 F = 99\,650.1420 \;>\; \log_2 p = 99\,648.5570. \checkmark
\]

\textbf{Step 3 — Witness $w_2 = 5$ for $q = 2$.}
\begin{align*}
  5^{p-1} &\equiv 1 \pmod{p}, \checkmark \\
  \gcd(5^{(p-1)/2}-1,\,p) &= 1. \checkmark
\end{align*}

\textbf{Step 4 — Witness $w_3 = 7$ for $q = 3$.}
\begin{align*}
  7^{p-1} &\equiv 1 \pmod{p}, \checkmark \\
  \gcd(7^{(p-1)/3}-1,\,p) &= 1. \checkmark
\end{align*}

\textbf{Wall-clock time:} $6\,945$ s $\approx 1$ h $56$ min
(single core, Python GMP, no multi-threading).\\[4pt]

\textbf{Conclusion:} all four Pocklington–Lehmer conditions are satisfied.
The number $p$ is \emph{unconditionally prime}.
\end{rigorousbox}

%------------------------------------------------------------------------------
\subsection{Comparison with the Original Computation}
%------------------------------------------------------------------------------

\begin{table}[h]
\centering
\caption{Original search vs.\ independent certificate verification
         for the 29\,998-digit prime.}
\label{tab:comparison29998}
\renewcommand{\arraystretch}{1.3}
\begin{tabular}{@{}lll@{}}
\toprule
 & \textbf{PrimeQuest v4 (search)} & \textbf{Independent verification} \\
\midrule
Purpose        & Find a prime candidate          & Check the certificate \\
Machine        & ARM64, 9 cores (Snapdragon X)   & Single core (Python GMP) \\
Wall-clock     & $11\,146$ s ($3$ h $06$ min)    & $6\,945$ s ($1$ h $56$ min) \\
MR tests run   & $286$ (probabilistic pre-filter) & — (no probabilistic step) \\
Result         & Certificate produced             & Certificate confirmed \\
\bottomrule
\end{tabular}
\end{table}

The verification is $\approx 1.6\times$ faster than the original search
because it executes exactly four modular exponentiations on the known
candidate, bypassing the sieve, Theorem-2 filter, and Miller–Rabin
pre-screening required during the search phase.

%------------------------------------------------------------------------------
\subsection{Record in Context}
%------------------------------------------------------------------------------

Table~\ref{tab:allcerts} (Section~\ref{sec:results}) lists all four
\textsc{PrimeQuest} certificates. The 29\,998-digit instance is currently
the largest unconditional Pocklington–Lehmer certificate in the family
$p = 3m(m+1)+1$, $m = 2^a\cdot 3^b - 1$, produced on a single consumer
laptop without distributed computing or specialised hardware.

%==============================================================================
% MINIMAL STAND-ALONE PREAMBLE (uncomment for self-contained compilation)
%==============================================================================
% \iffalse   % <-- remove this line to enable stand-alone mode
% \documentclass[11pt,a4paper]{article}
% \usepackage[utf8]{inputenc}
% \usepackage[T1]{fontenc}
% \usepackage[english]{babel}
% \usepackage{amsmath,amssymb,amsthm,mathtools,booktabs}
% \usepackage{geometry}\geometry{margin=2.5cm}
% \usepackage[dvipsnames]{xcolor}
% \usepackage{tcolorbox}
% \tcbuselibrary{skins,breakable}
% \newtcolorbox{rigorousbox}[1][]{colback=Green!8,colframe=Green!60!black,
%   arc=1mm,title=\textbf{Rigorous},fonttitle=\bfseries,breakable,#1}
% \newtcolorbox{resultbox}[1][]{colback=RoyalBlue!5,colframe=RoyalBlue!70!black,
%   arc=1mm,title=\textbf{Computational Result},fonttitle=\bfseries,breakable,#1}
% \newtheorem{theorem}{Theorem}[section]
% \newtheorem{lemma}[theorem]{Lemma}
% \theoremstyle{definition}
% \newtheorem{definition}[theorem]{Definition}
% \begin{document}
% %==============================================================================
% Unconditional Primality Certificate — 29 998-digit prime
% Parametric family  p = 3m(m+1)+1,  m = 2^a·3^b − 1
% Independent numerical verification — May 2026
%==============================================================================
%
% Stand-alone usage (requires paper_pocklington.tex preamble, or use the
% minimal preamble at the bottom of this file):
%
%   %==============================================================================
% Unconditional Primality Certificate — 29 998-digit prime
% Parametric family  p = 3m(m+1)+1,  m = 2^a·3^b − 1
% Independent numerical verification — May 2026
%==============================================================================
%
% Stand-alone usage (requires paper_pocklington.tex preamble, or use the
% minimal preamble at the bottom of this file):
%
%   %==============================================================================
% Unconditional Primality Certificate — 29 998-digit prime
% Parametric family  p = 3m(m+1)+1,  m = 2^a·3^b − 1
% Independent numerical verification — May 2026
%==============================================================================
%
% Stand-alone usage (requires paper_pocklington.tex preamble, or use the
% minimal preamble at the bottom of this file):
%
%   \input{certificate_29998}          % inside the main document
%
% Self-contained compilation:
%   pdflatex certificate_29998.tex
%
%==============================================================================

%------------------------------------------------------------------------------
\section{Unconditional Certificate for the 29\,998-Digit Prime}
\label{sec:cert29998}
%------------------------------------------------------------------------------

This section provides the complete, independently verified Pocklington–Lehmer
primality certificate for the largest prime currently produced by
\textsc{PrimeQuest}.

%------------------------------------------------------------------------------
\subsection{Parameters and Structure}
%------------------------------------------------------------------------------

\begin{definition}[29\,998-digit PrimeQuest instance]\label{def:instance29998}
Set
\[
  a = 19\,435,\qquad b = 19\,173,
\]
and define
\[
  m \;=\; 2^{19435}\cdot 3^{19173} - 1
    \qquad\bigl(14\,999 \text{ decimal digits}\bigr),
\]
\[
  p \;=\; 3m(m+1)+1
    \qquad\bigl(29\,998 \text{ decimal digits}\bigr).
\]
\end{definition}

The ratio $b/a = 19\,173/19\,435 \approx 0.9865$ is close to the
theoretical optimum
\[
  \rho^* = \frac{\log_{10} 2}{\log_{10} 3} \approx 0.6309
\]
in the sense that it minimises the digit count for a prescribed $a$;
the value used here targets exactly $29\,998$ digits and was located by
the \textsc{PrimeQuest v4} algorithm (see Section~\ref{sec:algo}).

\medskip

\noindent\textbf{Identification.}\quad
The first and last decimal digits of $p$ are:
\begin{equation}\label{eq:p29998digits}
  p \;=\; 1602087359509060348500037851549319056715\cdots
         16930852697403293697.
\end{equation}

%------------------------------------------------------------------------------
\subsection{Factorisation of $p-1$}
%------------------------------------------------------------------------------

\begin{rigorousbox}
\begin{lemma}[Factorisation of $p-1$]\label{lem:factorisation29998}
With $a, b, m, p$ as in Definition~\ref{def:instance29998},
\[
  p - 1 \;=\; \underbrace{2^{19435}\cdot 3^{19174}}_{F}
              \;\cdot\; m,
\]
where $F := 2^{19435}\cdot 3^{19174}$ has $15\,000$ decimal digits and
satisfies $F > \sqrt{p}$.
\end{lemma}
\end{rigorousbox}

\begin{proof}
Since $m + 1 = 2^{19435}\cdot 3^{19173}$,
\[
  p - 1 \;=\; 3m(m+1) \;=\; 3m\cdot 2^{19435}\cdot 3^{19173}
           \;=\; 2^{19435}\cdot 3^{19174}\cdot m.
\]
For the bound $F > \sqrt{p}$, note that $p < 3(m+1)^2$, so
$\sqrt{p} < \sqrt{3}\,(m+1) < 3(m+1) = 2^{19435}\cdot 3^{19174} = F$.

\noindent Numerically, comparing base-$2$ logarithms:
\[
  2\log_2 F
  \;=\; 2\bigl(19435 + 19174\log_2 3\bigr)
  \;\approx\; 99\,650.14,
\]
\[
  \log_2 p \;\approx\; 99\,648.56,
\]
confirming $F^2 > p$ with a margin of $\approx 1.58$ bits.
\end{proof}

%------------------------------------------------------------------------------
\subsection{The Pocklington--Lehmer Certificate}
%------------------------------------------------------------------------------

Theorem~\ref{thm:Pocklington} (recalled from Section~\ref{sec:family})
requires, for each prime $q \mid F$ — here $q \in \{2, 3\}$ — an integer
$w_q$ satisfying
\[
  w_q^{p-1} \equiv 1 \pmod{p}
  \qquad\text{and}\qquad
  \gcd\!\bigl(w_q^{(p-1)/q} - 1,\; p\bigr) = 1.
\]

\begin{resultbox}[title={\textbf{Certificate} — 29\,998-digit prime}]
\[
  p \;=\; 3m(m+1)+1,\qquad m = 2^{19435}\cdot 3^{19173}-1.
\]

\medskip
\noindent\textbf{Factored part of $p-1$:}
$\;F = 2^{19435}\cdot 3^{19174},\quad F > \sqrt{p}.\quad$\checkmark

\medskip
\renewcommand{\arraystretch}{1.4}
\begin{tabular}{@{}c l l@{}}
\toprule
Prime $q$ & Witness $w_q$ & Conditions verified \\
\midrule
$2$ & $w_2 = 5$ &
  $5^{p-1}\equiv 1\pmod{p}$\quad and\quad
  $\gcd(5^{(p-1)/2}-1,\,p)=1$ \\
$3$ & $w_3 = 7$ &
  $7^{p-1}\equiv 1\pmod{p}$\quad and\quad
  $\gcd(7^{(p-1)/3}-1,\,p)=1$ \\
\bottomrule
\end{tabular}

\medskip
\noindent\textbf{Conclusion:} $p$ is prime (deterministic proof). \checkmark
\end{resultbox}

%------------------------------------------------------------------------------
\subsection{Independent Numerical Verification}
%------------------------------------------------------------------------------

All four conditions of the certificate were verified independently using
\texttt{Python 3.12} with its built-in arbitrary-precision integer
arithmetic (GMP backend). The script performs no probabilistic step;
every operation is exact modular arithmetic.

\begin{rigorousbox}[title={\textbf{Verification record} — independent computation}]
\textbf{Environment:} Python 3.12, \texttt{sys.set\_int\_max\_str\_digits(40000)};
modular exponentiations via \texttt{pow(base, exp, mod)}.\\[4pt]

\textbf{Step 1 — Reconstructing $p$.}\quad
Computed $m = 2^{19435}\cdot 3^{19173}-1$ (14\,999 digits) and
$p = 3m(m+1)+1$ (29\,998 digits) from scratch.
First 40 digits of $p$: \texttt{1602087359509060348500037851549319056715};
last 20 digits: \texttt{16930852697403293697}.
Both match the stored certificate exactly.\\[4pt]

\textbf{Step 2 — Pocklington bound.}
\[
  2\log_2 F = 99\,650.1420 \;>\; \log_2 p = 99\,648.5570. \checkmark
\]

\textbf{Step 3 — Witness $w_2 = 5$ for $q = 2$.}
\begin{align*}
  5^{p-1} &\equiv 1 \pmod{p}, \checkmark \\
  \gcd(5^{(p-1)/2}-1,\,p) &= 1. \checkmark
\end{align*}

\textbf{Step 4 — Witness $w_3 = 7$ for $q = 3$.}
\begin{align*}
  7^{p-1} &\equiv 1 \pmod{p}, \checkmark \\
  \gcd(7^{(p-1)/3}-1,\,p) &= 1. \checkmark
\end{align*}

\textbf{Wall-clock time:} $6\,945$ s $\approx 1$ h $56$ min
(single core, Python GMP, no multi-threading).\\[4pt]

\textbf{Conclusion:} all four Pocklington–Lehmer conditions are satisfied.
The number $p$ is \emph{unconditionally prime}.
\end{rigorousbox}

%------------------------------------------------------------------------------
\subsection{Comparison with the Original Computation}
%------------------------------------------------------------------------------

\begin{table}[h]
\centering
\caption{Original search vs.\ independent certificate verification
         for the 29\,998-digit prime.}
\label{tab:comparison29998}
\renewcommand{\arraystretch}{1.3}
\begin{tabular}{@{}lll@{}}
\toprule
 & \textbf{PrimeQuest v4 (search)} & \textbf{Independent verification} \\
\midrule
Purpose        & Find a prime candidate          & Check the certificate \\
Machine        & ARM64, 9 cores (Snapdragon X)   & Single core (Python GMP) \\
Wall-clock     & $11\,146$ s ($3$ h $06$ min)    & $6\,945$ s ($1$ h $56$ min) \\
MR tests run   & $286$ (probabilistic pre-filter) & — (no probabilistic step) \\
Result         & Certificate produced             & Certificate confirmed \\
\bottomrule
\end{tabular}
\end{table}

The verification is $\approx 1.6\times$ faster than the original search
because it executes exactly four modular exponentiations on the known
candidate, bypassing the sieve, Theorem-2 filter, and Miller–Rabin
pre-screening required during the search phase.

%------------------------------------------------------------------------------
\subsection{Record in Context}
%------------------------------------------------------------------------------

Table~\ref{tab:allcerts} (Section~\ref{sec:results}) lists all four
\textsc{PrimeQuest} certificates. The 29\,998-digit instance is currently
the largest unconditional Pocklington–Lehmer certificate in the family
$p = 3m(m+1)+1$, $m = 2^a\cdot 3^b - 1$, produced on a single consumer
laptop without distributed computing or specialised hardware.

%==============================================================================
% MINIMAL STAND-ALONE PREAMBLE (uncomment for self-contained compilation)
%==============================================================================
% \iffalse   % <-- remove this line to enable stand-alone mode
% \documentclass[11pt,a4paper]{article}
% \usepackage[utf8]{inputenc}
% \usepackage[T1]{fontenc}
% \usepackage[english]{babel}
% \usepackage{amsmath,amssymb,amsthm,mathtools,booktabs}
% \usepackage{geometry}\geometry{margin=2.5cm}
% \usepackage[dvipsnames]{xcolor}
% \usepackage{tcolorbox}
% \tcbuselibrary{skins,breakable}
% \newtcolorbox{rigorousbox}[1][]{colback=Green!8,colframe=Green!60!black,
%   arc=1mm,title=\textbf{Rigorous},fonttitle=\bfseries,breakable,#1}
% \newtcolorbox{resultbox}[1][]{colback=RoyalBlue!5,colframe=RoyalBlue!70!black,
%   arc=1mm,title=\textbf{Computational Result},fonttitle=\bfseries,breakable,#1}
% \newtheorem{theorem}{Theorem}[section]
% \newtheorem{lemma}[theorem]{Lemma}
% \theoremstyle{definition}
% \newtheorem{definition}[theorem]{Definition}
% \begin{document}
% \input{certificate_29998}
% \end{document}
% \fi
          % inside the main document
%
% Self-contained compilation:
%   pdflatex certificate_29998.tex
%
%==============================================================================

%------------------------------------------------------------------------------
\section{Unconditional Certificate for the 29\,998-Digit Prime}
\label{sec:cert29998}
%------------------------------------------------------------------------------

This section provides the complete, independently verified Pocklington–Lehmer
primality certificate for the largest prime currently produced by
\textsc{PrimeQuest}.

%------------------------------------------------------------------------------
\subsection{Parameters and Structure}
%------------------------------------------------------------------------------

\begin{definition}[29\,998-digit PrimeQuest instance]\label{def:instance29998}
Set
\[
  a = 19\,435,\qquad b = 19\,173,
\]
and define
\[
  m \;=\; 2^{19435}\cdot 3^{19173} - 1
    \qquad\bigl(14\,999 \text{ decimal digits}\bigr),
\]
\[
  p \;=\; 3m(m+1)+1
    \qquad\bigl(29\,998 \text{ decimal digits}\bigr).
\]
\end{definition}

The ratio $b/a = 19\,173/19\,435 \approx 0.9865$ is close to the
theoretical optimum
\[
  \rho^* = \frac{\log_{10} 2}{\log_{10} 3} \approx 0.6309
\]
in the sense that it minimises the digit count for a prescribed $a$;
the value used here targets exactly $29\,998$ digits and was located by
the \textsc{PrimeQuest v4} algorithm (see Section~\ref{sec:algo}).

\medskip

\noindent\textbf{Identification.}\quad
The first and last decimal digits of $p$ are:
\begin{equation}\label{eq:p29998digits}
  p \;=\; 1602087359509060348500037851549319056715\cdots
         16930852697403293697.
\end{equation}

%------------------------------------------------------------------------------
\subsection{Factorisation of $p-1$}
%------------------------------------------------------------------------------

\begin{rigorousbox}
\begin{lemma}[Factorisation of $p-1$]\label{lem:factorisation29998}
With $a, b, m, p$ as in Definition~\ref{def:instance29998},
\[
  p - 1 \;=\; \underbrace{2^{19435}\cdot 3^{19174}}_{F}
              \;\cdot\; m,
\]
where $F := 2^{19435}\cdot 3^{19174}$ has $15\,000$ decimal digits and
satisfies $F > \sqrt{p}$.
\end{lemma}
\end{rigorousbox}

\begin{proof}
Since $m + 1 = 2^{19435}\cdot 3^{19173}$,
\[
  p - 1 \;=\; 3m(m+1) \;=\; 3m\cdot 2^{19435}\cdot 3^{19173}
           \;=\; 2^{19435}\cdot 3^{19174}\cdot m.
\]
For the bound $F > \sqrt{p}$, note that $p < 3(m+1)^2$, so
$\sqrt{p} < \sqrt{3}\,(m+1) < 3(m+1) = 2^{19435}\cdot 3^{19174} = F$.

\noindent Numerically, comparing base-$2$ logarithms:
\[
  2\log_2 F
  \;=\; 2\bigl(19435 + 19174\log_2 3\bigr)
  \;\approx\; 99\,650.14,
\]
\[
  \log_2 p \;\approx\; 99\,648.56,
\]
confirming $F^2 > p$ with a margin of $\approx 1.58$ bits.
\end{proof}

%------------------------------------------------------------------------------
\subsection{The Pocklington--Lehmer Certificate}
%------------------------------------------------------------------------------

Theorem~\ref{thm:Pocklington} (recalled from Section~\ref{sec:family})
requires, for each prime $q \mid F$ — here $q \in \{2, 3\}$ — an integer
$w_q$ satisfying
\[
  w_q^{p-1} \equiv 1 \pmod{p}
  \qquad\text{and}\qquad
  \gcd\!\bigl(w_q^{(p-1)/q} - 1,\; p\bigr) = 1.
\]

\begin{resultbox}[title={\textbf{Certificate} — 29\,998-digit prime}]
\[
  p \;=\; 3m(m+1)+1,\qquad m = 2^{19435}\cdot 3^{19173}-1.
\]

\medskip
\noindent\textbf{Factored part of $p-1$:}
$\;F = 2^{19435}\cdot 3^{19174},\quad F > \sqrt{p}.\quad$\checkmark

\medskip
\renewcommand{\arraystretch}{1.4}
\begin{tabular}{@{}c l l@{}}
\toprule
Prime $q$ & Witness $w_q$ & Conditions verified \\
\midrule
$2$ & $w_2 = 5$ &
  $5^{p-1}\equiv 1\pmod{p}$\quad and\quad
  $\gcd(5^{(p-1)/2}-1,\,p)=1$ \\
$3$ & $w_3 = 7$ &
  $7^{p-1}\equiv 1\pmod{p}$\quad and\quad
  $\gcd(7^{(p-1)/3}-1,\,p)=1$ \\
\bottomrule
\end{tabular}

\medskip
\noindent\textbf{Conclusion:} $p$ is prime (deterministic proof). \checkmark
\end{resultbox}

%------------------------------------------------------------------------------
\subsection{Independent Numerical Verification}
%------------------------------------------------------------------------------

All four conditions of the certificate were verified independently using
\texttt{Python 3.12} with its built-in arbitrary-precision integer
arithmetic (GMP backend). The script performs no probabilistic step;
every operation is exact modular arithmetic.

\begin{rigorousbox}[title={\textbf{Verification record} — independent computation}]
\textbf{Environment:} Python 3.12, \texttt{sys.set\_int\_max\_str\_digits(40000)};
modular exponentiations via \texttt{pow(base, exp, mod)}.\\[4pt]

\textbf{Step 1 — Reconstructing $p$.}\quad
Computed $m = 2^{19435}\cdot 3^{19173}-1$ (14\,999 digits) and
$p = 3m(m+1)+1$ (29\,998 digits) from scratch.
First 40 digits of $p$: \texttt{1602087359509060348500037851549319056715};
last 20 digits: \texttt{16930852697403293697}.
Both match the stored certificate exactly.\\[4pt]

\textbf{Step 2 — Pocklington bound.}
\[
  2\log_2 F = 99\,650.1420 \;>\; \log_2 p = 99\,648.5570. \checkmark
\]

\textbf{Step 3 — Witness $w_2 = 5$ for $q = 2$.}
\begin{align*}
  5^{p-1} &\equiv 1 \pmod{p}, \checkmark \\
  \gcd(5^{(p-1)/2}-1,\,p) &= 1. \checkmark
\end{align*}

\textbf{Step 4 — Witness $w_3 = 7$ for $q = 3$.}
\begin{align*}
  7^{p-1} &\equiv 1 \pmod{p}, \checkmark \\
  \gcd(7^{(p-1)/3}-1,\,p) &= 1. \checkmark
\end{align*}

\textbf{Wall-clock time:} $6\,945$ s $\approx 1$ h $56$ min
(single core, Python GMP, no multi-threading).\\[4pt]

\textbf{Conclusion:} all four Pocklington–Lehmer conditions are satisfied.
The number $p$ is \emph{unconditionally prime}.
\end{rigorousbox}

%------------------------------------------------------------------------------
\subsection{Comparison with the Original Computation}
%------------------------------------------------------------------------------

\begin{table}[h]
\centering
\caption{Original search vs.\ independent certificate verification
         for the 29\,998-digit prime.}
\label{tab:comparison29998}
\renewcommand{\arraystretch}{1.3}
\begin{tabular}{@{}lll@{}}
\toprule
 & \textbf{PrimeQuest v4 (search)} & \textbf{Independent verification} \\
\midrule
Purpose        & Find a prime candidate          & Check the certificate \\
Machine        & ARM64, 9 cores (Snapdragon X)   & Single core (Python GMP) \\
Wall-clock     & $11\,146$ s ($3$ h $06$ min)    & $6\,945$ s ($1$ h $56$ min) \\
MR tests run   & $286$ (probabilistic pre-filter) & — (no probabilistic step) \\
Result         & Certificate produced             & Certificate confirmed \\
\bottomrule
\end{tabular}
\end{table}

The verification is $\approx 1.6\times$ faster than the original search
because it executes exactly four modular exponentiations on the known
candidate, bypassing the sieve, Theorem-2 filter, and Miller–Rabin
pre-screening required during the search phase.

%------------------------------------------------------------------------------
\subsection{Record in Context}
%------------------------------------------------------------------------------

Table~\ref{tab:allcerts} (Section~\ref{sec:results}) lists all four
\textsc{PrimeQuest} certificates. The 29\,998-digit instance is currently
the largest unconditional Pocklington–Lehmer certificate in the family
$p = 3m(m+1)+1$, $m = 2^a\cdot 3^b - 1$, produced on a single consumer
laptop without distributed computing or specialised hardware.

%==============================================================================
% MINIMAL STAND-ALONE PREAMBLE (uncomment for self-contained compilation)
%==============================================================================
% \iffalse   % <-- remove this line to enable stand-alone mode
% \documentclass[11pt,a4paper]{article}
% \usepackage[utf8]{inputenc}
% \usepackage[T1]{fontenc}
% \usepackage[english]{babel}
% \usepackage{amsmath,amssymb,amsthm,mathtools,booktabs}
% \usepackage{geometry}\geometry{margin=2.5cm}
% \usepackage[dvipsnames]{xcolor}
% \usepackage{tcolorbox}
% \tcbuselibrary{skins,breakable}
% \newtcolorbox{rigorousbox}[1][]{colback=Green!8,colframe=Green!60!black,
%   arc=1mm,title=\textbf{Rigorous},fonttitle=\bfseries,breakable,#1}
% \newtcolorbox{resultbox}[1][]{colback=RoyalBlue!5,colframe=RoyalBlue!70!black,
%   arc=1mm,title=\textbf{Computational Result},fonttitle=\bfseries,breakable,#1}
% \newtheorem{theorem}{Theorem}[section]
% \newtheorem{lemma}[theorem]{Lemma}
% \theoremstyle{definition}
% \newtheorem{definition}[theorem]{Definition}
% \begin{document}
% %==============================================================================
% Unconditional Primality Certificate — 29 998-digit prime
% Parametric family  p = 3m(m+1)+1,  m = 2^a·3^b − 1
% Independent numerical verification — May 2026
%==============================================================================
%
% Stand-alone usage (requires paper_pocklington.tex preamble, or use the
% minimal preamble at the bottom of this file):
%
%   \input{certificate_29998}          % inside the main document
%
% Self-contained compilation:
%   pdflatex certificate_29998.tex
%
%==============================================================================

%------------------------------------------------------------------------------
\section{Unconditional Certificate for the 29\,998-Digit Prime}
\label{sec:cert29998}
%------------------------------------------------------------------------------

This section provides the complete, independently verified Pocklington–Lehmer
primality certificate for the largest prime currently produced by
\textsc{PrimeQuest}.

%------------------------------------------------------------------------------
\subsection{Parameters and Structure}
%------------------------------------------------------------------------------

\begin{definition}[29\,998-digit PrimeQuest instance]\label{def:instance29998}
Set
\[
  a = 19\,435,\qquad b = 19\,173,
\]
and define
\[
  m \;=\; 2^{19435}\cdot 3^{19173} - 1
    \qquad\bigl(14\,999 \text{ decimal digits}\bigr),
\]
\[
  p \;=\; 3m(m+1)+1
    \qquad\bigl(29\,998 \text{ decimal digits}\bigr).
\]
\end{definition}

The ratio $b/a = 19\,173/19\,435 \approx 0.9865$ is close to the
theoretical optimum
\[
  \rho^* = \frac{\log_{10} 2}{\log_{10} 3} \approx 0.6309
\]
in the sense that it minimises the digit count for a prescribed $a$;
the value used here targets exactly $29\,998$ digits and was located by
the \textsc{PrimeQuest v4} algorithm (see Section~\ref{sec:algo}).

\medskip

\noindent\textbf{Identification.}\quad
The first and last decimal digits of $p$ are:
\begin{equation}\label{eq:p29998digits}
  p \;=\; 1602087359509060348500037851549319056715\cdots
         16930852697403293697.
\end{equation}

%------------------------------------------------------------------------------
\subsection{Factorisation of $p-1$}
%------------------------------------------------------------------------------

\begin{rigorousbox}
\begin{lemma}[Factorisation of $p-1$]\label{lem:factorisation29998}
With $a, b, m, p$ as in Definition~\ref{def:instance29998},
\[
  p - 1 \;=\; \underbrace{2^{19435}\cdot 3^{19174}}_{F}
              \;\cdot\; m,
\]
where $F := 2^{19435}\cdot 3^{19174}$ has $15\,000$ decimal digits and
satisfies $F > \sqrt{p}$.
\end{lemma}
\end{rigorousbox}

\begin{proof}
Since $m + 1 = 2^{19435}\cdot 3^{19173}$,
\[
  p - 1 \;=\; 3m(m+1) \;=\; 3m\cdot 2^{19435}\cdot 3^{19173}
           \;=\; 2^{19435}\cdot 3^{19174}\cdot m.
\]
For the bound $F > \sqrt{p}$, note that $p < 3(m+1)^2$, so
$\sqrt{p} < \sqrt{3}\,(m+1) < 3(m+1) = 2^{19435}\cdot 3^{19174} = F$.

\noindent Numerically, comparing base-$2$ logarithms:
\[
  2\log_2 F
  \;=\; 2\bigl(19435 + 19174\log_2 3\bigr)
  \;\approx\; 99\,650.14,
\]
\[
  \log_2 p \;\approx\; 99\,648.56,
\]
confirming $F^2 > p$ with a margin of $\approx 1.58$ bits.
\end{proof}

%------------------------------------------------------------------------------
\subsection{The Pocklington--Lehmer Certificate}
%------------------------------------------------------------------------------

Theorem~\ref{thm:Pocklington} (recalled from Section~\ref{sec:family})
requires, for each prime $q \mid F$ — here $q \in \{2, 3\}$ — an integer
$w_q$ satisfying
\[
  w_q^{p-1} \equiv 1 \pmod{p}
  \qquad\text{and}\qquad
  \gcd\!\bigl(w_q^{(p-1)/q} - 1,\; p\bigr) = 1.
\]

\begin{resultbox}[title={\textbf{Certificate} — 29\,998-digit prime}]
\[
  p \;=\; 3m(m+1)+1,\qquad m = 2^{19435}\cdot 3^{19173}-1.
\]

\medskip
\noindent\textbf{Factored part of $p-1$:}
$\;F = 2^{19435}\cdot 3^{19174},\quad F > \sqrt{p}.\quad$\checkmark

\medskip
\renewcommand{\arraystretch}{1.4}
\begin{tabular}{@{}c l l@{}}
\toprule
Prime $q$ & Witness $w_q$ & Conditions verified \\
\midrule
$2$ & $w_2 = 5$ &
  $5^{p-1}\equiv 1\pmod{p}$\quad and\quad
  $\gcd(5^{(p-1)/2}-1,\,p)=1$ \\
$3$ & $w_3 = 7$ &
  $7^{p-1}\equiv 1\pmod{p}$\quad and\quad
  $\gcd(7^{(p-1)/3}-1,\,p)=1$ \\
\bottomrule
\end{tabular}

\medskip
\noindent\textbf{Conclusion:} $p$ is prime (deterministic proof). \checkmark
\end{resultbox}

%------------------------------------------------------------------------------
\subsection{Independent Numerical Verification}
%------------------------------------------------------------------------------

All four conditions of the certificate were verified independently using
\texttt{Python 3.12} with its built-in arbitrary-precision integer
arithmetic (GMP backend). The script performs no probabilistic step;
every operation is exact modular arithmetic.

\begin{rigorousbox}[title={\textbf{Verification record} — independent computation}]
\textbf{Environment:} Python 3.12, \texttt{sys.set\_int\_max\_str\_digits(40000)};
modular exponentiations via \texttt{pow(base, exp, mod)}.\\[4pt]

\textbf{Step 1 — Reconstructing $p$.}\quad
Computed $m = 2^{19435}\cdot 3^{19173}-1$ (14\,999 digits) and
$p = 3m(m+1)+1$ (29\,998 digits) from scratch.
First 40 digits of $p$: \texttt{1602087359509060348500037851549319056715};
last 20 digits: \texttt{16930852697403293697}.
Both match the stored certificate exactly.\\[4pt]

\textbf{Step 2 — Pocklington bound.}
\[
  2\log_2 F = 99\,650.1420 \;>\; \log_2 p = 99\,648.5570. \checkmark
\]

\textbf{Step 3 — Witness $w_2 = 5$ for $q = 2$.}
\begin{align*}
  5^{p-1} &\equiv 1 \pmod{p}, \checkmark \\
  \gcd(5^{(p-1)/2}-1,\,p) &= 1. \checkmark
\end{align*}

\textbf{Step 4 — Witness $w_3 = 7$ for $q = 3$.}
\begin{align*}
  7^{p-1} &\equiv 1 \pmod{p}, \checkmark \\
  \gcd(7^{(p-1)/3}-1,\,p) &= 1. \checkmark
\end{align*}

\textbf{Wall-clock time:} $6\,945$ s $\approx 1$ h $56$ min
(single core, Python GMP, no multi-threading).\\[4pt]

\textbf{Conclusion:} all four Pocklington–Lehmer conditions are satisfied.
The number $p$ is \emph{unconditionally prime}.
\end{rigorousbox}

%------------------------------------------------------------------------------
\subsection{Comparison with the Original Computation}
%------------------------------------------------------------------------------

\begin{table}[h]
\centering
\caption{Original search vs.\ independent certificate verification
         for the 29\,998-digit prime.}
\label{tab:comparison29998}
\renewcommand{\arraystretch}{1.3}
\begin{tabular}{@{}lll@{}}
\toprule
 & \textbf{PrimeQuest v4 (search)} & \textbf{Independent verification} \\
\midrule
Purpose        & Find a prime candidate          & Check the certificate \\
Machine        & ARM64, 9 cores (Snapdragon X)   & Single core (Python GMP) \\
Wall-clock     & $11\,146$ s ($3$ h $06$ min)    & $6\,945$ s ($1$ h $56$ min) \\
MR tests run   & $286$ (probabilistic pre-filter) & — (no probabilistic step) \\
Result         & Certificate produced             & Certificate confirmed \\
\bottomrule
\end{tabular}
\end{table}

The verification is $\approx 1.6\times$ faster than the original search
because it executes exactly four modular exponentiations on the known
candidate, bypassing the sieve, Theorem-2 filter, and Miller–Rabin
pre-screening required during the search phase.

%------------------------------------------------------------------------------
\subsection{Record in Context}
%------------------------------------------------------------------------------

Table~\ref{tab:allcerts} (Section~\ref{sec:results}) lists all four
\textsc{PrimeQuest} certificates. The 29\,998-digit instance is currently
the largest unconditional Pocklington–Lehmer certificate in the family
$p = 3m(m+1)+1$, $m = 2^a\cdot 3^b - 1$, produced on a single consumer
laptop without distributed computing or specialised hardware.

%==============================================================================
% MINIMAL STAND-ALONE PREAMBLE (uncomment for self-contained compilation)
%==============================================================================
% \iffalse   % <-- remove this line to enable stand-alone mode
% \documentclass[11pt,a4paper]{article}
% \usepackage[utf8]{inputenc}
% \usepackage[T1]{fontenc}
% \usepackage[english]{babel}
% \usepackage{amsmath,amssymb,amsthm,mathtools,booktabs}
% \usepackage{geometry}\geometry{margin=2.5cm}
% \usepackage[dvipsnames]{xcolor}
% \usepackage{tcolorbox}
% \tcbuselibrary{skins,breakable}
% \newtcolorbox{rigorousbox}[1][]{colback=Green!8,colframe=Green!60!black,
%   arc=1mm,title=\textbf{Rigorous},fonttitle=\bfseries,breakable,#1}
% \newtcolorbox{resultbox}[1][]{colback=RoyalBlue!5,colframe=RoyalBlue!70!black,
%   arc=1mm,title=\textbf{Computational Result},fonttitle=\bfseries,breakable,#1}
% \newtheorem{theorem}{Theorem}[section]
% \newtheorem{lemma}[theorem]{Lemma}
% \theoremstyle{definition}
% \newtheorem{definition}[theorem]{Definition}
% \begin{document}
% \input{certificate_29998}
% \end{document}
% \fi

% \end{document}
% \fi
          % inside the main document
%
% Self-contained compilation:
%   pdflatex certificate_29998.tex
%
%==============================================================================

%------------------------------------------------------------------------------
\section{Unconditional Certificate for the 29\,998-Digit Prime}
\label{sec:cert29998}
%------------------------------------------------------------------------------

This section provides the complete, independently verified Pocklington–Lehmer
primality certificate for the largest prime currently produced by
\textsc{PrimeQuest}.

%------------------------------------------------------------------------------
\subsection{Parameters and Structure}
%------------------------------------------------------------------------------

\begin{definition}[29\,998-digit PrimeQuest instance]\label{def:instance29998}
Set
\[
  a = 19\,435,\qquad b = 19\,173,
\]
and define
\[
  m \;=\; 2^{19435}\cdot 3^{19173} - 1
    \qquad\bigl(14\,999 \text{ decimal digits}\bigr),
\]
\[
  p \;=\; 3m(m+1)+1
    \qquad\bigl(29\,998 \text{ decimal digits}\bigr).
\]
\end{definition}

The ratio $b/a = 19\,173/19\,435 \approx 0.9865$ is close to the
theoretical optimum
\[
  \rho^* = \frac{\log_{10} 2}{\log_{10} 3} \approx 0.6309
\]
in the sense that it minimises the digit count for a prescribed $a$;
the value used here targets exactly $29\,998$ digits and was located by
the \textsc{PrimeQuest v4} algorithm (see Section~\ref{sec:algo}).

\medskip

\noindent\textbf{Identification.}\quad
The first and last decimal digits of $p$ are:
\begin{equation}\label{eq:p29998digits}
  p \;=\; 1602087359509060348500037851549319056715\cdots
         16930852697403293697.
\end{equation}

%------------------------------------------------------------------------------
\subsection{Factorisation of $p-1$}
%------------------------------------------------------------------------------

\begin{rigorousbox}
\begin{lemma}[Factorisation of $p-1$]\label{lem:factorisation29998}
With $a, b, m, p$ as in Definition~\ref{def:instance29998},
\[
  p - 1 \;=\; \underbrace{2^{19435}\cdot 3^{19174}}_{F}
              \;\cdot\; m,
\]
where $F := 2^{19435}\cdot 3^{19174}$ has $15\,000$ decimal digits and
satisfies $F > \sqrt{p}$.
\end{lemma}
\end{rigorousbox}

\begin{proof}
Since $m + 1 = 2^{19435}\cdot 3^{19173}$,
\[
  p - 1 \;=\; 3m(m+1) \;=\; 3m\cdot 2^{19435}\cdot 3^{19173}
           \;=\; 2^{19435}\cdot 3^{19174}\cdot m.
\]
For the bound $F > \sqrt{p}$, note that $p < 3(m+1)^2$, so
$\sqrt{p} < \sqrt{3}\,(m+1) < 3(m+1) = 2^{19435}\cdot 3^{19174} = F$.

\noindent Numerically, comparing base-$2$ logarithms:
\[
  2\log_2 F
  \;=\; 2\bigl(19435 + 19174\log_2 3\bigr)
  \;\approx\; 99\,650.14,
\]
\[
  \log_2 p \;\approx\; 99\,648.56,
\]
confirming $F^2 > p$ with a margin of $\approx 1.58$ bits.
\end{proof}

%------------------------------------------------------------------------------
\subsection{The Pocklington--Lehmer Certificate}
%------------------------------------------------------------------------------

Theorem~\ref{thm:Pocklington} (recalled from Section~\ref{sec:family})
requires, for each prime $q \mid F$ — here $q \in \{2, 3\}$ — an integer
$w_q$ satisfying
\[
  w_q^{p-1} \equiv 1 \pmod{p}
  \qquad\text{and}\qquad
  \gcd\!\bigl(w_q^{(p-1)/q} - 1,\; p\bigr) = 1.
\]

\begin{resultbox}[title={\textbf{Certificate} — 29\,998-digit prime}]
\[
  p \;=\; 3m(m+1)+1,\qquad m = 2^{19435}\cdot 3^{19173}-1.
\]

\medskip
\noindent\textbf{Factored part of $p-1$:}
$\;F = 2^{19435}\cdot 3^{19174},\quad F > \sqrt{p}.\quad$\checkmark

\medskip
\renewcommand{\arraystretch}{1.4}
\begin{tabular}{@{}c l l@{}}
\toprule
Prime $q$ & Witness $w_q$ & Conditions verified \\
\midrule
$2$ & $w_2 = 5$ &
  $5^{p-1}\equiv 1\pmod{p}$\quad and\quad
  $\gcd(5^{(p-1)/2}-1,\,p)=1$ \\
$3$ & $w_3 = 7$ &
  $7^{p-1}\equiv 1\pmod{p}$\quad and\quad
  $\gcd(7^{(p-1)/3}-1,\,p)=1$ \\
\bottomrule
\end{tabular}

\medskip
\noindent\textbf{Conclusion:} $p$ is prime (deterministic proof). \checkmark
\end{resultbox}

%------------------------------------------------------------------------------
\subsection{Independent Numerical Verification}
%------------------------------------------------------------------------------

All four conditions of the certificate were verified independently using
\texttt{Python 3.12} with its built-in arbitrary-precision integer
arithmetic (GMP backend). The script performs no probabilistic step;
every operation is exact modular arithmetic.

\begin{rigorousbox}[title={\textbf{Verification record} — independent computation}]
\textbf{Environment:} Python 3.12, \texttt{sys.set\_int\_max\_str\_digits(40000)};
modular exponentiations via \texttt{pow(base, exp, mod)}.\\[4pt]

\textbf{Step 1 — Reconstructing $p$.}\quad
Computed $m = 2^{19435}\cdot 3^{19173}-1$ (14\,999 digits) and
$p = 3m(m+1)+1$ (29\,998 digits) from scratch.
First 40 digits of $p$: \texttt{1602087359509060348500037851549319056715};
last 20 digits: \texttt{16930852697403293697}.
Both match the stored certificate exactly.\\[4pt]

\textbf{Step 2 — Pocklington bound.}
\[
  2\log_2 F = 99\,650.1420 \;>\; \log_2 p = 99\,648.5570. \checkmark
\]

\textbf{Step 3 — Witness $w_2 = 5$ for $q = 2$.}
\begin{align*}
  5^{p-1} &\equiv 1 \pmod{p}, \checkmark \\
  \gcd(5^{(p-1)/2}-1,\,p) &= 1. \checkmark
\end{align*}

\textbf{Step 4 — Witness $w_3 = 7$ for $q = 3$.}
\begin{align*}
  7^{p-1} &\equiv 1 \pmod{p}, \checkmark \\
  \gcd(7^{(p-1)/3}-1,\,p) &= 1. \checkmark
\end{align*}

\textbf{Wall-clock time:} $6\,945$ s $\approx 1$ h $56$ min
(single core, Python GMP, no multi-threading).\\[4pt]

\textbf{Conclusion:} all four Pocklington–Lehmer conditions are satisfied.
The number $p$ is \emph{unconditionally prime}.
\end{rigorousbox}

%------------------------------------------------------------------------------
\subsection{Comparison with the Original Computation}
%------------------------------------------------------------------------------

\begin{table}[h]
\centering
\caption{Original search vs.\ independent certificate verification
         for the 29\,998-digit prime.}
\label{tab:comparison29998}
\renewcommand{\arraystretch}{1.3}
\begin{tabular}{@{}lll@{}}
\toprule
 & \textbf{PrimeQuest v4 (search)} & \textbf{Independent verification} \\
\midrule
Purpose        & Find a prime candidate          & Check the certificate \\
Machine        & ARM64, 9 cores (Snapdragon X)   & Single core (Python GMP) \\
Wall-clock     & $11\,146$ s ($3$ h $06$ min)    & $6\,945$ s ($1$ h $56$ min) \\
MR tests run   & $286$ (probabilistic pre-filter) & — (no probabilistic step) \\
Result         & Certificate produced             & Certificate confirmed \\
\bottomrule
\end{tabular}
\end{table}

The verification is $\approx 1.6\times$ faster than the original search
because it executes exactly four modular exponentiations on the known
candidate, bypassing the sieve, Theorem-2 filter, and Miller–Rabin
pre-screening required during the search phase.

%------------------------------------------------------------------------------
\subsection{Record in Context}
%------------------------------------------------------------------------------

Table~\ref{tab:allcerts} (Section~\ref{sec:results}) lists all four
\textsc{PrimeQuest} certificates. The 29\,998-digit instance is currently
the largest unconditional Pocklington–Lehmer certificate in the family
$p = 3m(m+1)+1$, $m = 2^a\cdot 3^b - 1$, produced on a single consumer
laptop without distributed computing or specialised hardware.

%==============================================================================
% MINIMAL STAND-ALONE PREAMBLE (uncomment for self-contained compilation)
%==============================================================================
% \iffalse   % <-- remove this line to enable stand-alone mode
% \documentclass[11pt,a4paper]{article}
% \usepackage[utf8]{inputenc}
% \usepackage[T1]{fontenc}
% \usepackage[english]{babel}
% \usepackage{amsmath,amssymb,amsthm,mathtools,booktabs}
% \usepackage{geometry}\geometry{margin=2.5cm}
% \usepackage[dvipsnames]{xcolor}
% \usepackage{tcolorbox}
% \tcbuselibrary{skins,breakable}
% \newtcolorbox{rigorousbox}[1][]{colback=Green!8,colframe=Green!60!black,
%   arc=1mm,title=\textbf{Rigorous},fonttitle=\bfseries,breakable,#1}
% \newtcolorbox{resultbox}[1][]{colback=RoyalBlue!5,colframe=RoyalBlue!70!black,
%   arc=1mm,title=\textbf{Computational Result},fonttitle=\bfseries,breakable,#1}
% \newtheorem{theorem}{Theorem}[section]
% \newtheorem{lemma}[theorem]{Lemma}
% \theoremstyle{definition}
% \newtheorem{definition}[theorem]{Definition}
% \begin{document}
% %==============================================================================
% Unconditional Primality Certificate — 29 998-digit prime
% Parametric family  p = 3m(m+1)+1,  m = 2^a·3^b − 1
% Independent numerical verification — May 2026
%==============================================================================
%
% Stand-alone usage (requires paper_pocklington.tex preamble, or use the
% minimal preamble at the bottom of this file):
%
%   %==============================================================================
% Unconditional Primality Certificate — 29 998-digit prime
% Parametric family  p = 3m(m+1)+1,  m = 2^a·3^b − 1
% Independent numerical verification — May 2026
%==============================================================================
%
% Stand-alone usage (requires paper_pocklington.tex preamble, or use the
% minimal preamble at the bottom of this file):
%
%   \input{certificate_29998}          % inside the main document
%
% Self-contained compilation:
%   pdflatex certificate_29998.tex
%
%==============================================================================

%------------------------------------------------------------------------------
\section{Unconditional Certificate for the 29\,998-Digit Prime}
\label{sec:cert29998}
%------------------------------------------------------------------------------

This section provides the complete, independently verified Pocklington–Lehmer
primality certificate for the largest prime currently produced by
\textsc{PrimeQuest}.

%------------------------------------------------------------------------------
\subsection{Parameters and Structure}
%------------------------------------------------------------------------------

\begin{definition}[29\,998-digit PrimeQuest instance]\label{def:instance29998}
Set
\[
  a = 19\,435,\qquad b = 19\,173,
\]
and define
\[
  m \;=\; 2^{19435}\cdot 3^{19173} - 1
    \qquad\bigl(14\,999 \text{ decimal digits}\bigr),
\]
\[
  p \;=\; 3m(m+1)+1
    \qquad\bigl(29\,998 \text{ decimal digits}\bigr).
\]
\end{definition}

The ratio $b/a = 19\,173/19\,435 \approx 0.9865$ is close to the
theoretical optimum
\[
  \rho^* = \frac{\log_{10} 2}{\log_{10} 3} \approx 0.6309
\]
in the sense that it minimises the digit count for a prescribed $a$;
the value used here targets exactly $29\,998$ digits and was located by
the \textsc{PrimeQuest v4} algorithm (see Section~\ref{sec:algo}).

\medskip

\noindent\textbf{Identification.}\quad
The first and last decimal digits of $p$ are:
\begin{equation}\label{eq:p29998digits}
  p \;=\; 1602087359509060348500037851549319056715\cdots
         16930852697403293697.
\end{equation}

%------------------------------------------------------------------------------
\subsection{Factorisation of $p-1$}
%------------------------------------------------------------------------------

\begin{rigorousbox}
\begin{lemma}[Factorisation of $p-1$]\label{lem:factorisation29998}
With $a, b, m, p$ as in Definition~\ref{def:instance29998},
\[
  p - 1 \;=\; \underbrace{2^{19435}\cdot 3^{19174}}_{F}
              \;\cdot\; m,
\]
where $F := 2^{19435}\cdot 3^{19174}$ has $15\,000$ decimal digits and
satisfies $F > \sqrt{p}$.
\end{lemma}
\end{rigorousbox}

\begin{proof}
Since $m + 1 = 2^{19435}\cdot 3^{19173}$,
\[
  p - 1 \;=\; 3m(m+1) \;=\; 3m\cdot 2^{19435}\cdot 3^{19173}
           \;=\; 2^{19435}\cdot 3^{19174}\cdot m.
\]
For the bound $F > \sqrt{p}$, note that $p < 3(m+1)^2$, so
$\sqrt{p} < \sqrt{3}\,(m+1) < 3(m+1) = 2^{19435}\cdot 3^{19174} = F$.

\noindent Numerically, comparing base-$2$ logarithms:
\[
  2\log_2 F
  \;=\; 2\bigl(19435 + 19174\log_2 3\bigr)
  \;\approx\; 99\,650.14,
\]
\[
  \log_2 p \;\approx\; 99\,648.56,
\]
confirming $F^2 > p$ with a margin of $\approx 1.58$ bits.
\end{proof}

%------------------------------------------------------------------------------
\subsection{The Pocklington--Lehmer Certificate}
%------------------------------------------------------------------------------

Theorem~\ref{thm:Pocklington} (recalled from Section~\ref{sec:family})
requires, for each prime $q \mid F$ — here $q \in \{2, 3\}$ — an integer
$w_q$ satisfying
\[
  w_q^{p-1} \equiv 1 \pmod{p}
  \qquad\text{and}\qquad
  \gcd\!\bigl(w_q^{(p-1)/q} - 1,\; p\bigr) = 1.
\]

\begin{resultbox}[title={\textbf{Certificate} — 29\,998-digit prime}]
\[
  p \;=\; 3m(m+1)+1,\qquad m = 2^{19435}\cdot 3^{19173}-1.
\]

\medskip
\noindent\textbf{Factored part of $p-1$:}
$\;F = 2^{19435}\cdot 3^{19174},\quad F > \sqrt{p}.\quad$\checkmark

\medskip
\renewcommand{\arraystretch}{1.4}
\begin{tabular}{@{}c l l@{}}
\toprule
Prime $q$ & Witness $w_q$ & Conditions verified \\
\midrule
$2$ & $w_2 = 5$ &
  $5^{p-1}\equiv 1\pmod{p}$\quad and\quad
  $\gcd(5^{(p-1)/2}-1,\,p)=1$ \\
$3$ & $w_3 = 7$ &
  $7^{p-1}\equiv 1\pmod{p}$\quad and\quad
  $\gcd(7^{(p-1)/3}-1,\,p)=1$ \\
\bottomrule
\end{tabular}

\medskip
\noindent\textbf{Conclusion:} $p$ is prime (deterministic proof). \checkmark
\end{resultbox}

%------------------------------------------------------------------------------
\subsection{Independent Numerical Verification}
%------------------------------------------------------------------------------

All four conditions of the certificate were verified independently using
\texttt{Python 3.12} with its built-in arbitrary-precision integer
arithmetic (GMP backend). The script performs no probabilistic step;
every operation is exact modular arithmetic.

\begin{rigorousbox}[title={\textbf{Verification record} — independent computation}]
\textbf{Environment:} Python 3.12, \texttt{sys.set\_int\_max\_str\_digits(40000)};
modular exponentiations via \texttt{pow(base, exp, mod)}.\\[4pt]

\textbf{Step 1 — Reconstructing $p$.}\quad
Computed $m = 2^{19435}\cdot 3^{19173}-1$ (14\,999 digits) and
$p = 3m(m+1)+1$ (29\,998 digits) from scratch.
First 40 digits of $p$: \texttt{1602087359509060348500037851549319056715};
last 20 digits: \texttt{16930852697403293697}.
Both match the stored certificate exactly.\\[4pt]

\textbf{Step 2 — Pocklington bound.}
\[
  2\log_2 F = 99\,650.1420 \;>\; \log_2 p = 99\,648.5570. \checkmark
\]

\textbf{Step 3 — Witness $w_2 = 5$ for $q = 2$.}
\begin{align*}
  5^{p-1} &\equiv 1 \pmod{p}, \checkmark \\
  \gcd(5^{(p-1)/2}-1,\,p) &= 1. \checkmark
\end{align*}

\textbf{Step 4 — Witness $w_3 = 7$ for $q = 3$.}
\begin{align*}
  7^{p-1} &\equiv 1 \pmod{p}, \checkmark \\
  \gcd(7^{(p-1)/3}-1,\,p) &= 1. \checkmark
\end{align*}

\textbf{Wall-clock time:} $6\,945$ s $\approx 1$ h $56$ min
(single core, Python GMP, no multi-threading).\\[4pt]

\textbf{Conclusion:} all four Pocklington–Lehmer conditions are satisfied.
The number $p$ is \emph{unconditionally prime}.
\end{rigorousbox}

%------------------------------------------------------------------------------
\subsection{Comparison with the Original Computation}
%------------------------------------------------------------------------------

\begin{table}[h]
\centering
\caption{Original search vs.\ independent certificate verification
         for the 29\,998-digit prime.}
\label{tab:comparison29998}
\renewcommand{\arraystretch}{1.3}
\begin{tabular}{@{}lll@{}}
\toprule
 & \textbf{PrimeQuest v4 (search)} & \textbf{Independent verification} \\
\midrule
Purpose        & Find a prime candidate          & Check the certificate \\
Machine        & ARM64, 9 cores (Snapdragon X)   & Single core (Python GMP) \\
Wall-clock     & $11\,146$ s ($3$ h $06$ min)    & $6\,945$ s ($1$ h $56$ min) \\
MR tests run   & $286$ (probabilistic pre-filter) & — (no probabilistic step) \\
Result         & Certificate produced             & Certificate confirmed \\
\bottomrule
\end{tabular}
\end{table}

The verification is $\approx 1.6\times$ faster than the original search
because it executes exactly four modular exponentiations on the known
candidate, bypassing the sieve, Theorem-2 filter, and Miller–Rabin
pre-screening required during the search phase.

%------------------------------------------------------------------------------
\subsection{Record in Context}
%------------------------------------------------------------------------------

Table~\ref{tab:allcerts} (Section~\ref{sec:results}) lists all four
\textsc{PrimeQuest} certificates. The 29\,998-digit instance is currently
the largest unconditional Pocklington–Lehmer certificate in the family
$p = 3m(m+1)+1$, $m = 2^a\cdot 3^b - 1$, produced on a single consumer
laptop without distributed computing or specialised hardware.

%==============================================================================
% MINIMAL STAND-ALONE PREAMBLE (uncomment for self-contained compilation)
%==============================================================================
% \iffalse   % <-- remove this line to enable stand-alone mode
% \documentclass[11pt,a4paper]{article}
% \usepackage[utf8]{inputenc}
% \usepackage[T1]{fontenc}
% \usepackage[english]{babel}
% \usepackage{amsmath,amssymb,amsthm,mathtools,booktabs}
% \usepackage{geometry}\geometry{margin=2.5cm}
% \usepackage[dvipsnames]{xcolor}
% \usepackage{tcolorbox}
% \tcbuselibrary{skins,breakable}
% \newtcolorbox{rigorousbox}[1][]{colback=Green!8,colframe=Green!60!black,
%   arc=1mm,title=\textbf{Rigorous},fonttitle=\bfseries,breakable,#1}
% \newtcolorbox{resultbox}[1][]{colback=RoyalBlue!5,colframe=RoyalBlue!70!black,
%   arc=1mm,title=\textbf{Computational Result},fonttitle=\bfseries,breakable,#1}
% \newtheorem{theorem}{Theorem}[section]
% \newtheorem{lemma}[theorem]{Lemma}
% \theoremstyle{definition}
% \newtheorem{definition}[theorem]{Definition}
% \begin{document}
% \input{certificate_29998}
% \end{document}
% \fi
          % inside the main document
%
% Self-contained compilation:
%   pdflatex certificate_29998.tex
%
%==============================================================================

%------------------------------------------------------------------------------
\section{Unconditional Certificate for the 29\,998-Digit Prime}
\label{sec:cert29998}
%------------------------------------------------------------------------------

This section provides the complete, independently verified Pocklington–Lehmer
primality certificate for the largest prime currently produced by
\textsc{PrimeQuest}.

%------------------------------------------------------------------------------
\subsection{Parameters and Structure}
%------------------------------------------------------------------------------

\begin{definition}[29\,998-digit PrimeQuest instance]\label{def:instance29998}
Set
\[
  a = 19\,435,\qquad b = 19\,173,
\]
and define
\[
  m \;=\; 2^{19435}\cdot 3^{19173} - 1
    \qquad\bigl(14\,999 \text{ decimal digits}\bigr),
\]
\[
  p \;=\; 3m(m+1)+1
    \qquad\bigl(29\,998 \text{ decimal digits}\bigr).
\]
\end{definition}

The ratio $b/a = 19\,173/19\,435 \approx 0.9865$ is close to the
theoretical optimum
\[
  \rho^* = \frac{\log_{10} 2}{\log_{10} 3} \approx 0.6309
\]
in the sense that it minimises the digit count for a prescribed $a$;
the value used here targets exactly $29\,998$ digits and was located by
the \textsc{PrimeQuest v4} algorithm (see Section~\ref{sec:algo}).

\medskip

\noindent\textbf{Identification.}\quad
The first and last decimal digits of $p$ are:
\begin{equation}\label{eq:p29998digits}
  p \;=\; 1602087359509060348500037851549319056715\cdots
         16930852697403293697.
\end{equation}

%------------------------------------------------------------------------------
\subsection{Factorisation of $p-1$}
%------------------------------------------------------------------------------

\begin{rigorousbox}
\begin{lemma}[Factorisation of $p-1$]\label{lem:factorisation29998}
With $a, b, m, p$ as in Definition~\ref{def:instance29998},
\[
  p - 1 \;=\; \underbrace{2^{19435}\cdot 3^{19174}}_{F}
              \;\cdot\; m,
\]
where $F := 2^{19435}\cdot 3^{19174}$ has $15\,000$ decimal digits and
satisfies $F > \sqrt{p}$.
\end{lemma}
\end{rigorousbox}

\begin{proof}
Since $m + 1 = 2^{19435}\cdot 3^{19173}$,
\[
  p - 1 \;=\; 3m(m+1) \;=\; 3m\cdot 2^{19435}\cdot 3^{19173}
           \;=\; 2^{19435}\cdot 3^{19174}\cdot m.
\]
For the bound $F > \sqrt{p}$, note that $p < 3(m+1)^2$, so
$\sqrt{p} < \sqrt{3}\,(m+1) < 3(m+1) = 2^{19435}\cdot 3^{19174} = F$.

\noindent Numerically, comparing base-$2$ logarithms:
\[
  2\log_2 F
  \;=\; 2\bigl(19435 + 19174\log_2 3\bigr)
  \;\approx\; 99\,650.14,
\]
\[
  \log_2 p \;\approx\; 99\,648.56,
\]
confirming $F^2 > p$ with a margin of $\approx 1.58$ bits.
\end{proof}

%------------------------------------------------------------------------------
\subsection{The Pocklington--Lehmer Certificate}
%------------------------------------------------------------------------------

Theorem~\ref{thm:Pocklington} (recalled from Section~\ref{sec:family})
requires, for each prime $q \mid F$ — here $q \in \{2, 3\}$ — an integer
$w_q$ satisfying
\[
  w_q^{p-1} \equiv 1 \pmod{p}
  \qquad\text{and}\qquad
  \gcd\!\bigl(w_q^{(p-1)/q} - 1,\; p\bigr) = 1.
\]

\begin{resultbox}[title={\textbf{Certificate} — 29\,998-digit prime}]
\[
  p \;=\; 3m(m+1)+1,\qquad m = 2^{19435}\cdot 3^{19173}-1.
\]

\medskip
\noindent\textbf{Factored part of $p-1$:}
$\;F = 2^{19435}\cdot 3^{19174},\quad F > \sqrt{p}.\quad$\checkmark

\medskip
\renewcommand{\arraystretch}{1.4}
\begin{tabular}{@{}c l l@{}}
\toprule
Prime $q$ & Witness $w_q$ & Conditions verified \\
\midrule
$2$ & $w_2 = 5$ &
  $5^{p-1}\equiv 1\pmod{p}$\quad and\quad
  $\gcd(5^{(p-1)/2}-1,\,p)=1$ \\
$3$ & $w_3 = 7$ &
  $7^{p-1}\equiv 1\pmod{p}$\quad and\quad
  $\gcd(7^{(p-1)/3}-1,\,p)=1$ \\
\bottomrule
\end{tabular}

\medskip
\noindent\textbf{Conclusion:} $p$ is prime (deterministic proof). \checkmark
\end{resultbox}

%------------------------------------------------------------------------------
\subsection{Independent Numerical Verification}
%------------------------------------------------------------------------------

All four conditions of the certificate were verified independently using
\texttt{Python 3.12} with its built-in arbitrary-precision integer
arithmetic (GMP backend). The script performs no probabilistic step;
every operation is exact modular arithmetic.

\begin{rigorousbox}[title={\textbf{Verification record} — independent computation}]
\textbf{Environment:} Python 3.12, \texttt{sys.set\_int\_max\_str\_digits(40000)};
modular exponentiations via \texttt{pow(base, exp, mod)}.\\[4pt]

\textbf{Step 1 — Reconstructing $p$.}\quad
Computed $m = 2^{19435}\cdot 3^{19173}-1$ (14\,999 digits) and
$p = 3m(m+1)+1$ (29\,998 digits) from scratch.
First 40 digits of $p$: \texttt{1602087359509060348500037851549319056715};
last 20 digits: \texttt{16930852697403293697}.
Both match the stored certificate exactly.\\[4pt]

\textbf{Step 2 — Pocklington bound.}
\[
  2\log_2 F = 99\,650.1420 \;>\; \log_2 p = 99\,648.5570. \checkmark
\]

\textbf{Step 3 — Witness $w_2 = 5$ for $q = 2$.}
\begin{align*}
  5^{p-1} &\equiv 1 \pmod{p}, \checkmark \\
  \gcd(5^{(p-1)/2}-1,\,p) &= 1. \checkmark
\end{align*}

\textbf{Step 4 — Witness $w_3 = 7$ for $q = 3$.}
\begin{align*}
  7^{p-1} &\equiv 1 \pmod{p}, \checkmark \\
  \gcd(7^{(p-1)/3}-1,\,p) &= 1. \checkmark
\end{align*}

\textbf{Wall-clock time:} $6\,945$ s $\approx 1$ h $56$ min
(single core, Python GMP, no multi-threading).\\[4pt]

\textbf{Conclusion:} all four Pocklington–Lehmer conditions are satisfied.
The number $p$ is \emph{unconditionally prime}.
\end{rigorousbox}

%------------------------------------------------------------------------------
\subsection{Comparison with the Original Computation}
%------------------------------------------------------------------------------

\begin{table}[h]
\centering
\caption{Original search vs.\ independent certificate verification
         for the 29\,998-digit prime.}
\label{tab:comparison29998}
\renewcommand{\arraystretch}{1.3}
\begin{tabular}{@{}lll@{}}
\toprule
 & \textbf{PrimeQuest v4 (search)} & \textbf{Independent verification} \\
\midrule
Purpose        & Find a prime candidate          & Check the certificate \\
Machine        & ARM64, 9 cores (Snapdragon X)   & Single core (Python GMP) \\
Wall-clock     & $11\,146$ s ($3$ h $06$ min)    & $6\,945$ s ($1$ h $56$ min) \\
MR tests run   & $286$ (probabilistic pre-filter) & — (no probabilistic step) \\
Result         & Certificate produced             & Certificate confirmed \\
\bottomrule
\end{tabular}
\end{table}

The verification is $\approx 1.6\times$ faster than the original search
because it executes exactly four modular exponentiations on the known
candidate, bypassing the sieve, Theorem-2 filter, and Miller–Rabin
pre-screening required during the search phase.

%------------------------------------------------------------------------------
\subsection{Record in Context}
%------------------------------------------------------------------------------

Table~\ref{tab:allcerts} (Section~\ref{sec:results}) lists all four
\textsc{PrimeQuest} certificates. The 29\,998-digit instance is currently
the largest unconditional Pocklington–Lehmer certificate in the family
$p = 3m(m+1)+1$, $m = 2^a\cdot 3^b - 1$, produced on a single consumer
laptop without distributed computing or specialised hardware.

%==============================================================================
% MINIMAL STAND-ALONE PREAMBLE (uncomment for self-contained compilation)
%==============================================================================
% \iffalse   % <-- remove this line to enable stand-alone mode
% \documentclass[11pt,a4paper]{article}
% \usepackage[utf8]{inputenc}
% \usepackage[T1]{fontenc}
% \usepackage[english]{babel}
% \usepackage{amsmath,amssymb,amsthm,mathtools,booktabs}
% \usepackage{geometry}\geometry{margin=2.5cm}
% \usepackage[dvipsnames]{xcolor}
% \usepackage{tcolorbox}
% \tcbuselibrary{skins,breakable}
% \newtcolorbox{rigorousbox}[1][]{colback=Green!8,colframe=Green!60!black,
%   arc=1mm,title=\textbf{Rigorous},fonttitle=\bfseries,breakable,#1}
% \newtcolorbox{resultbox}[1][]{colback=RoyalBlue!5,colframe=RoyalBlue!70!black,
%   arc=1mm,title=\textbf{Computational Result},fonttitle=\bfseries,breakable,#1}
% \newtheorem{theorem}{Theorem}[section]
% \newtheorem{lemma}[theorem]{Lemma}
% \theoremstyle{definition}
% \newtheorem{definition}[theorem]{Definition}
% \begin{document}
% %==============================================================================
% Unconditional Primality Certificate — 29 998-digit prime
% Parametric family  p = 3m(m+1)+1,  m = 2^a·3^b − 1
% Independent numerical verification — May 2026
%==============================================================================
%
% Stand-alone usage (requires paper_pocklington.tex preamble, or use the
% minimal preamble at the bottom of this file):
%
%   \input{certificate_29998}          % inside the main document
%
% Self-contained compilation:
%   pdflatex certificate_29998.tex
%
%==============================================================================

%------------------------------------------------------------------------------
\section{Unconditional Certificate for the 29\,998-Digit Prime}
\label{sec:cert29998}
%------------------------------------------------------------------------------

This section provides the complete, independently verified Pocklington–Lehmer
primality certificate for the largest prime currently produced by
\textsc{PrimeQuest}.

%------------------------------------------------------------------------------
\subsection{Parameters and Structure}
%------------------------------------------------------------------------------

\begin{definition}[29\,998-digit PrimeQuest instance]\label{def:instance29998}
Set
\[
  a = 19\,435,\qquad b = 19\,173,
\]
and define
\[
  m \;=\; 2^{19435}\cdot 3^{19173} - 1
    \qquad\bigl(14\,999 \text{ decimal digits}\bigr),
\]
\[
  p \;=\; 3m(m+1)+1
    \qquad\bigl(29\,998 \text{ decimal digits}\bigr).
\]
\end{definition}

The ratio $b/a = 19\,173/19\,435 \approx 0.9865$ is close to the
theoretical optimum
\[
  \rho^* = \frac{\log_{10} 2}{\log_{10} 3} \approx 0.6309
\]
in the sense that it minimises the digit count for a prescribed $a$;
the value used here targets exactly $29\,998$ digits and was located by
the \textsc{PrimeQuest v4} algorithm (see Section~\ref{sec:algo}).

\medskip

\noindent\textbf{Identification.}\quad
The first and last decimal digits of $p$ are:
\begin{equation}\label{eq:p29998digits}
  p \;=\; 1602087359509060348500037851549319056715\cdots
         16930852697403293697.
\end{equation}

%------------------------------------------------------------------------------
\subsection{Factorisation of $p-1$}
%------------------------------------------------------------------------------

\begin{rigorousbox}
\begin{lemma}[Factorisation of $p-1$]\label{lem:factorisation29998}
With $a, b, m, p$ as in Definition~\ref{def:instance29998},
\[
  p - 1 \;=\; \underbrace{2^{19435}\cdot 3^{19174}}_{F}
              \;\cdot\; m,
\]
where $F := 2^{19435}\cdot 3^{19174}$ has $15\,000$ decimal digits and
satisfies $F > \sqrt{p}$.
\end{lemma}
\end{rigorousbox}

\begin{proof}
Since $m + 1 = 2^{19435}\cdot 3^{19173}$,
\[
  p - 1 \;=\; 3m(m+1) \;=\; 3m\cdot 2^{19435}\cdot 3^{19173}
           \;=\; 2^{19435}\cdot 3^{19174}\cdot m.
\]
For the bound $F > \sqrt{p}$, note that $p < 3(m+1)^2$, so
$\sqrt{p} < \sqrt{3}\,(m+1) < 3(m+1) = 2^{19435}\cdot 3^{19174} = F$.

\noindent Numerically, comparing base-$2$ logarithms:
\[
  2\log_2 F
  \;=\; 2\bigl(19435 + 19174\log_2 3\bigr)
  \;\approx\; 99\,650.14,
\]
\[
  \log_2 p \;\approx\; 99\,648.56,
\]
confirming $F^2 > p$ with a margin of $\approx 1.58$ bits.
\end{proof}

%------------------------------------------------------------------------------
\subsection{The Pocklington--Lehmer Certificate}
%------------------------------------------------------------------------------

Theorem~\ref{thm:Pocklington} (recalled from Section~\ref{sec:family})
requires, for each prime $q \mid F$ — here $q \in \{2, 3\}$ — an integer
$w_q$ satisfying
\[
  w_q^{p-1} \equiv 1 \pmod{p}
  \qquad\text{and}\qquad
  \gcd\!\bigl(w_q^{(p-1)/q} - 1,\; p\bigr) = 1.
\]

\begin{resultbox}[title={\textbf{Certificate} — 29\,998-digit prime}]
\[
  p \;=\; 3m(m+1)+1,\qquad m = 2^{19435}\cdot 3^{19173}-1.
\]

\medskip
\noindent\textbf{Factored part of $p-1$:}
$\;F = 2^{19435}\cdot 3^{19174},\quad F > \sqrt{p}.\quad$\checkmark

\medskip
\renewcommand{\arraystretch}{1.4}
\begin{tabular}{@{}c l l@{}}
\toprule
Prime $q$ & Witness $w_q$ & Conditions verified \\
\midrule
$2$ & $w_2 = 5$ &
  $5^{p-1}\equiv 1\pmod{p}$\quad and\quad
  $\gcd(5^{(p-1)/2}-1,\,p)=1$ \\
$3$ & $w_3 = 7$ &
  $7^{p-1}\equiv 1\pmod{p}$\quad and\quad
  $\gcd(7^{(p-1)/3}-1,\,p)=1$ \\
\bottomrule
\end{tabular}

\medskip
\noindent\textbf{Conclusion:} $p$ is prime (deterministic proof). \checkmark
\end{resultbox}

%------------------------------------------------------------------------------
\subsection{Independent Numerical Verification}
%------------------------------------------------------------------------------

All four conditions of the certificate were verified independently using
\texttt{Python 3.12} with its built-in arbitrary-precision integer
arithmetic (GMP backend). The script performs no probabilistic step;
every operation is exact modular arithmetic.

\begin{rigorousbox}[title={\textbf{Verification record} — independent computation}]
\textbf{Environment:} Python 3.12, \texttt{sys.set\_int\_max\_str\_digits(40000)};
modular exponentiations via \texttt{pow(base, exp, mod)}.\\[4pt]

\textbf{Step 1 — Reconstructing $p$.}\quad
Computed $m = 2^{19435}\cdot 3^{19173}-1$ (14\,999 digits) and
$p = 3m(m+1)+1$ (29\,998 digits) from scratch.
First 40 digits of $p$: \texttt{1602087359509060348500037851549319056715};
last 20 digits: \texttt{16930852697403293697}.
Both match the stored certificate exactly.\\[4pt]

\textbf{Step 2 — Pocklington bound.}
\[
  2\log_2 F = 99\,650.1420 \;>\; \log_2 p = 99\,648.5570. \checkmark
\]

\textbf{Step 3 — Witness $w_2 = 5$ for $q = 2$.}
\begin{align*}
  5^{p-1} &\equiv 1 \pmod{p}, \checkmark \\
  \gcd(5^{(p-1)/2}-1,\,p) &= 1. \checkmark
\end{align*}

\textbf{Step 4 — Witness $w_3 = 7$ for $q = 3$.}
\begin{align*}
  7^{p-1} &\equiv 1 \pmod{p}, \checkmark \\
  \gcd(7^{(p-1)/3}-1,\,p) &= 1. \checkmark
\end{align*}

\textbf{Wall-clock time:} $6\,945$ s $\approx 1$ h $56$ min
(single core, Python GMP, no multi-threading).\\[4pt]

\textbf{Conclusion:} all four Pocklington–Lehmer conditions are satisfied.
The number $p$ is \emph{unconditionally prime}.
\end{rigorousbox}

%------------------------------------------------------------------------------
\subsection{Comparison with the Original Computation}
%------------------------------------------------------------------------------

\begin{table}[h]
\centering
\caption{Original search vs.\ independent certificate verification
         for the 29\,998-digit prime.}
\label{tab:comparison29998}
\renewcommand{\arraystretch}{1.3}
\begin{tabular}{@{}lll@{}}
\toprule
 & \textbf{PrimeQuest v4 (search)} & \textbf{Independent verification} \\
\midrule
Purpose        & Find a prime candidate          & Check the certificate \\
Machine        & ARM64, 9 cores (Snapdragon X)   & Single core (Python GMP) \\
Wall-clock     & $11\,146$ s ($3$ h $06$ min)    & $6\,945$ s ($1$ h $56$ min) \\
MR tests run   & $286$ (probabilistic pre-filter) & — (no probabilistic step) \\
Result         & Certificate produced             & Certificate confirmed \\
\bottomrule
\end{tabular}
\end{table}

The verification is $\approx 1.6\times$ faster than the original search
because it executes exactly four modular exponentiations on the known
candidate, bypassing the sieve, Theorem-2 filter, and Miller–Rabin
pre-screening required during the search phase.

%------------------------------------------------------------------------------
\subsection{Record in Context}
%------------------------------------------------------------------------------

Table~\ref{tab:allcerts} (Section~\ref{sec:results}) lists all four
\textsc{PrimeQuest} certificates. The 29\,998-digit instance is currently
the largest unconditional Pocklington–Lehmer certificate in the family
$p = 3m(m+1)+1$, $m = 2^a\cdot 3^b - 1$, produced on a single consumer
laptop without distributed computing or specialised hardware.

%==============================================================================
% MINIMAL STAND-ALONE PREAMBLE (uncomment for self-contained compilation)
%==============================================================================
% \iffalse   % <-- remove this line to enable stand-alone mode
% \documentclass[11pt,a4paper]{article}
% \usepackage[utf8]{inputenc}
% \usepackage[T1]{fontenc}
% \usepackage[english]{babel}
% \usepackage{amsmath,amssymb,amsthm,mathtools,booktabs}
% \usepackage{geometry}\geometry{margin=2.5cm}
% \usepackage[dvipsnames]{xcolor}
% \usepackage{tcolorbox}
% \tcbuselibrary{skins,breakable}
% \newtcolorbox{rigorousbox}[1][]{colback=Green!8,colframe=Green!60!black,
%   arc=1mm,title=\textbf{Rigorous},fonttitle=\bfseries,breakable,#1}
% \newtcolorbox{resultbox}[1][]{colback=RoyalBlue!5,colframe=RoyalBlue!70!black,
%   arc=1mm,title=\textbf{Computational Result},fonttitle=\bfseries,breakable,#1}
% \newtheorem{theorem}{Theorem}[section]
% \newtheorem{lemma}[theorem]{Lemma}
% \theoremstyle{definition}
% \newtheorem{definition}[theorem]{Definition}
% \begin{document}
% \input{certificate_29998}
% \end{document}
% \fi

% \end{document}
% \fi

% \end{document}
% \fi

% \end{document}
% \fi
